% 原书全部参考文献，保留原语言与次序；条目边界和规范化审计见 references_work。
\chapter*{参考文献}
\addcontentsline{toc}{chapter}{参考文献}
\markboth{参考文献}{参考文献}
\begingroup
\small\linespread{1.15}\selectfont
\raggedright\raggedbottom
\setlength{\parskip}{0pt}
\clubpenalty=10000\widowpenalty=10000
\hyphenpenalty=10000\exhyphenpenalty=200
\begin{list}{}{\setlength{\leftmargin}{2.1em}\setlength{\labelwidth}{0pt}
  \setlength{\labelsep}{0pt}\setlength{\itemindent}{-2.1em}
  \setlength{\parsep}{0pt}\setlength{\itemsep}{2pt}\setlength{\topsep}{0pt}}
% REF001; 原书分章页 1; RP01-L002,RP01-L003
\item\label{refentry:REF001}O. Aalen, Nonparametric estimation of partial transition probabilities in multiple decrement models. Ann. Stat. \textbf{6}, 534--545 (1978)\label{refend:REF001}
% REF002; 原书分章页 1; RP01-L004,RP01-L005
\item\label{refentry:REF002}A. Abadie, G. Imbens, Simple and bias-corrected matching estimators for average treatment effects. Technical Report 283. NBER Working Paper (2002)\label{refend:REF002}
% REF003; 原书分章页 1; RP01-L006,RP01-L007
\item\label{refentry:REF003}Action to Control Cardiovascular Risk in Diabetes Study Group. Effects of intensive glucose lowering in type 2 diabetes. N. Engl. J. Med. \textbf{358}, 2545--2549 (2008)\label{refend:REF003}
% REF004; 原书分章页 1; RP01-L008,RP01-L009
\item\label{refentry:REF004}ADVANCE Collaborative Group, Intensive blood glucose control and vascular outcomes in patients with type 2 diabetes. N. Engl. J. Med. \textbf{358}, 2560--2562 (2008)\label{refend:REF004}
% REF005; 原书分章页 1; RP01-L010,RP01-L011
\item\label{refentry:REF005}A. Afifi, S. Azen, \textit{Statistical} \textit{Analysis:} \textit{A} \textit{Computer} \textit{Oriented} \textit{Approach}, 2nd edn. (Academic, New York, 1979)\label{refend:REF005}
% REF006; 原书分章页 1; RP01-L012,RP01-L013
\item\label{refentry:REF006}C. Anderson, The end of theory: the data deluge makes the scientific method obsolete. Wired (2008)\label{refend:REF006}
% REF007; 原书分章页 1; RP01-L014,RP01-L015,RP01-L016,RP01-L017,RP01-L018,RP01-L019,RP01-L020,RP01-L021
\item\label{refentry:REF007}T.J. Aragon, epitools: Epidemiology tools (2012). \url{http://cran.r-project.org/package=epitools}\label{refend:REF007}
% REF008; 原书分章页 1; RP01-L022,RP01-L023
\item\label{refentry:REF008}S. Aral, D. Walker, Identifying social influence in networks using randomized experiments. IEEE Intell. Syst. \textbf{26}(5), 91--96 (2011)\label{refend:REF008}
% REF009; 原书分章页 1; RP01-L024,RP01-L025
\item\label{refentry:REF009}S. Aral, D. Walker, Tie strength, embeddedness, and social influence: a large-scale networked experiment. Manag. Sci. \textbf{60}(6), 1352--1370 (2014)\label{refend:REF009}
% REF010; 原书分章页 1; RP01-L026,RP01-L027
\item\label{refentry:REF010}P.M. Aronow, C. Samii, Estimating average causal effects under interference between units. ArXiv e-prints, May (2013)\label{refend:REF010}
% REF011; 原书分章页 1; RP01-L028,RP01-L029
\item\label{refentry:REF011}J.Y. Audibert, A.B. Tsybakov, Fast learning rates for plug-in classifiers. Ann. Stat. \textbf{35}(2), 608--633 (2007)\label{refend:REF011}
% REF012; 原书分章页 1; RP01-L030,RP01-L031
\item\label{refentry:REF012}L. Auret, C. Aldrich, Empirical comparison of tree ensemble variable importance measures. Chemom. Intel. Lab. Syst. \textbf{105}(2), 157--170 (2011)\label{refend:REF012}
% REF013; 原书分章页 1; RP01-L032,RP01-L033,RP01-L034,RP01-L035
\item\label{refentry:REF013}P.C. Austin, A. Manca, M. Zwarensteina, D.N. Juurlinka, M.B. Stanbrook, A substantial and confusing variation exists in handling of baseline covariates in randomized controlled trials: a review of trials published in leading medical journals. J. Clin. Epidemiol. \textbf{63}, 142--153 (2010)\label{refend:REF013}
% REF014; 原书分章页 1; RP01-L036,RP01-L037
\item\label{refentry:REF014}C. Avin, I. Shpitser, J. Pearl, Identifiability of path-specific effects. Proceedings of International Joint Conference on Artificial Intelligence, 357--363 (2005)\label{refend:REF014}
% REF015; 原书分章页 2; RP02-L003,RP02-L004
\item\label{refentry:REF015}S. Balakrishnan, D. Madigan, Algorithms for sparse linear classifiers in the massive data setting. J. Mach. Learn. Res. \textbf{9}, 313--337 (2008)\label{refend:REF015}
% REF016; 原书分章页 2; RP02-L005,RP02-L006,RP02-L007
\item\label{refentry:REF016}L. Balzer, M. Petersen, M.J. van der Laan, Adaptive pair-matching in randomized trials with unbiased and efficient effect estimation. Stat. Med. \textbf{34}(6), 999--1011 (2015)\label{refend:REF016}
% REF017; 原书分章页 2; RP02-L008,RP02-L009,RP02-L010
\item\label{refentry:REF017}L. Balzer, J. Ahern, S. Galea, M.J. van der Laan, Estimating effects with rare outcomes and high dimensional covariates: Knowledge is power. Epidemiol. Methods. \textbf{5}(1), 1--18 (2016a)\label{refend:REF017}
% REF018; 原书分章页 2; RP02-L011,RP02-L012,RP02-L013
\item\label{refentry:REF018}L. Balzer, M. van der Laan, M. Petersen, the SEARCH Collaboration, Adaptive pre-specification in randomized trials with and without pair-matching. Stat. Med. \textbf{35}(25), 4528--4545 (2016b)\label{refend:REF018}
% REF019; 原书分章页 2; RP02-L014,RP02-L015,RP02-L016
\item\label{refentry:REF019}L.B. Balzer, M.L. Petersen, M.J. van der Laan, the SEARCH Collaboration, Targeted estimation and inference of the sample average treatment effect in trials with and without pair-matching. Stat. Med. \textbf{35}(21), 3717--3732 (2016c)\label{refend:REF019}
% REF020; 原书分章页 2; RP02-L017,RP02-L018
\item\label{refentry:REF020}H. Bang, J.M. Robins, Doubly robust estimation in missing data and causal inference models. Biometrics \textbf{61}, 962--972 (2005)\label{refend:REF020}
% REF021; 原书分章页 2; RP02-L019,RP02-L020
\item\label{refentry:REF021}A.-L. Barabási, R. Albert, Emergence of scaling in random networks. Science \textbf{286}(5439), 509--512 (1999)\label{refend:REF021}
% REF022; 原书分章页 2; RP02-L021,RP02-L022
\item\label{refentry:REF022}E. Bareinboim, J. Pearl, A general algorithm for deciding transportability of experimental results. J. Causal Inf. \textbf{1}(1), 107--134 (2013)\label{refend:REF022}
% REF023; 原书分章页 2; RP02-L023,RP02-L024
\item\label{refentry:REF023}G.W. Basse, E.M. Airoldi, Optimal design of experiments in the presence of network-correlated outcomes. ArXiv e-prints, July (2015)\label{refend:REF023}
% REF024; 原书分章页 2; RP02-L025,RP02-L026
\item\label{refentry:REF024}C. Beck, B. Lu, R. Greevy, nbpMatching: functions for optimal non-bipartite optimal matching (2016). \url{https://CRAN.R-project.org/package=nbpMatching}\label{refend:REF024}
% REF025; 原书分章页 2; RP02-L027,RP02-L028,RP02-L029
\item\label{refentry:REF025}O. Bembom, M.J. van der Laan, A practical illustration of the importance of realistic individualized treatment rules in causal inference. Electron. J. Stat. \textbf{1}, 574--596 (2007)\label{refend:REF025}
% REF026; 原书分章页 2; RP02-L030,RP02-L031,RP02-L032
\item\label{refentry:REF026}O. Bembom, M.J. van der Laan, Analyzing sequentially randomized trials based on causal effect models for realistic individualized treatment rules. Stat. Med. \textbf{27}, 3689--3716 (2008)\label{refend:REF026}
% REF027; 原书分章页 2; RP02-L033,RP02-L034,RP02-L035,RP02-L036
\item\label{refentry:REF027}O. Bembom, M.L. Petersen, S.-Y. Rhee, W.J. Fessel, S.E. Sinisi, R.W. Shafer, M.J. van der Laan, Biomarker discovery using targeted maximum likelihood estimation: application to the treatment of antiretroviral resistant HIV infection. Stat. Med. \textbf{28}, 152--72 (2009)\label{refend:REF027}
% REF028; 原书分章页 2; RP02-L037,RP02-L038
\item\label{refentry:REF028}J. Benichou, M.H. Gail, Estimates of absolute cause-specific risk in cohort studies. Biometrics \textbf{46}, 813--826 (1990)\label{refend:REF028}
% REF029; 原书分章页 2; RP02-L039,RP02-L040
\item\label{refentry:REF029}Y. Benjamini, Y. Hochberg, Controlling the false discovery rate: a practical and powerful approach to multiple testing. J. R. Stat. Soc. Ser. B \textbf{57}, 289--300 (1995)\label{refend:REF029}
% REF030; 原书分章页 2; RP02-L041,RP02-L042,RP02-L043
\item\label{refentry:REF030}D. Benkeser, M.J. van der Laan, The highly adaptive lasso estimator, in \textit{IEEE} \textit{International} \textit{Conference} \textit{on} \textit{Data} \textit{Science} \textit{and} \textit{Advanced} \textit{Analytics}, pp. 689--696 (2016)\label{refend:REF030}
% REF031; 原书分章页 2; RP02-L044,RP02-L045,RP02-L046
\item\label{refentry:REF031}D. Benkeser, M. Carone, M.J. van der Laan, P. Gilbert, Doubly-robust nonparametric inference on the average treatment effect. Biometrika. \textbf{104}(4), 863--880 (2017a)\label{refend:REF031}
% REF032; 原书分章页 3; RP03-L003,RP03-L004
\item\label{refentry:REF032}D. Benkeser, S.D. Lendle, J. Cheng, M.J. van der Laan, Online cross-validation-based ensemble learning. Stat. Med. \textbf{37}(2), 249--260 (2017b)\label{refend:REF032}
% REF033; 原书分章页 3; RP03-L005,RP03-L006
\item\label{refentry:REF033}P. Bertail, A. Chambaz, E. Joly, Practical targeted learning from large data sets by survey sampling. ArXiv e-prints, June (2016)\label{refend:REF033}
% REF034; 原书分章页 3; RP03-L007,RP03-L008
\item\label{refentry:REF034}P. Bertail, E. Chautru, S. Clémençon, Empirical processes in survey sampling with (conditional) Poisson designs. Scand. J. Stat. \textbf{44}(1), 97--111 (2017)\label{refend:REF034}
% REF035; 原书分章页 3; RP03-L009
\item\label{refentry:REF035}P.J. Bickel, On adaptive estimation. Ann. Stat. \textbf{10}, 647--671 (1982)\label{refend:REF035}
% REF036; 原书分章页 3; RP03-L010,RP03-L011
\item\label{refentry:REF036}P.J. Bickel, F. Götze, W.R. van Zwet, Resampling fewer than n observations: gains, losses, and remedies for losses. Stat. Sin. \textbf{7}(1), 1--31 (1997a)\label{refend:REF036}
% REF037; 原书分章页 3; RP03-L012,RP03-L013
\item\label{refentry:REF037}P.J. Bickel, C.A.J. Klaassen, Y. Ritov, J. Wellner, \textit{Efficient} \textit{and} \textit{Adaptive} \textit{Estimation} \textit{for} \textit{Semiparametric} \textit{Models} (Springer, Berlin, Heidelberg, New York, 1997b)\label{refend:REF037}
% REF038; 原书分章页 3; RP03-L014,RP03-L015
\item\label{refentry:REF038}L. Bondesson, I. Traat, A. Lundqvist, Pareto sampling versus Sampford and conditional Poisson sampling. Scand. J. Stat. Theory Appl. \textbf{33}(4), 699--720 (2006)\label{refend:REF038}
% REF039; 原书分章页 3; RP03-L016,RP03-L017
\item\label{refentry:REF039}L. Bottou, Large-scale machine learning with stochastic gradient descent, in \textit{Proceedings} \textit{of} \textit{COMPSTAT'2010} (Springer, Berlin, 2010), pp. 177--186\label{refend:REF039}
% REF040; 原书分章页 3; RP03-L018,RP03-L019
\item\label{refentry:REF040}L. Bottou, Stochastic gradient descent tricks, in \textit{Neural} \textit{Networks:} \textit{Tricks} \textit{of} \textit{the} \textit{Trade} (Springer, Berlin, 2012), pp. 421--436\label{refend:REF040}
% REF041; 原书分章页 3; RP03-L020,RP03-L021
\item\label{refentry:REF041}J. Bowers, M.M. Fredrickson, C. Panagopoulos, Reasoning about interference between units: a general framework. Polit. Anal. \textbf{21}(1), 97--124 (2013)\label{refend:REF041}
% REF042; 原书分章页 3; RP03-L022
\item\label{refentry:REF042}L. Breiman, Random forests. Mach. Learn. \textbf{45}, 5--32 (2001)\label{refend:REF042}
% REF043; 原书分章页 3; RP03-L023,RP03-L024
\item\label{refentry:REF043}L. Breiman, J.H. Friedman, R. Olshen, C.J. Stone, \textit{Classification} \textit{and} \textit{Regression} \textit{Trees} (Chapman \& Hall, Boca Raton, 1984)\label{refend:REF043}
% REF044; 原书分章页 3; RP03-L025,RP03-L026
\item\label{refentry:REF044}L. Breiman et al., Statistical modeling: the two cultures (with comments and a rejoinder by the author). Stat. Sci. \textbf{16}(3), 199--231 (2001)\label{refend:REF044}
% REF045; 原书分章页 3; RP03-L027,RP03-L028
\item\label{refentry:REF045}D.I. Broadhurst, D.B. Kell, Statistical strategies for avoiding false discoveries in metabolomics and related experiments. Metabolomics \textbf{2}(4), 171--196 (2006)\label{refend:REF045}
% REF046; 原书分章页 3; RP03-L029,RP03-L030
\item\label{refentry:REF046}D.W. Brock, D. Wikler, Ethical challenges in long-term funding for HIV/AIDS. Health Aff. \textbf{28}(6), 1666--1676 (2009)\label{refend:REF046}
% REF047; 原书分章页 3; RP03-L031,RP03-L032,RP03-L033
\item\label{refentry:REF047}J.C. Brooks, Super learner and targeted maximum likelihood estimation for longitudinal data structures with applications to atrial fibrillation. PhD thesis, University of California, Berkeley (2012)\label{refend:REF047}
% REF048; 原书分章页 3; RP03-L034,RP03-L035,RP03-L036
\item\label{refentry:REF048}L.E. Cain, J.M. Robins, E. Lanoy, R. Logan, D. Costagliola, M.A. Hernan, When to start treatment? A systematic approach to the comparison of dynamic regimes using observational data. Int. J. Biostat. \textbf{6}, Article 18 (2010)\label{refend:REF048}
% REF049; 原书分章页 3; RP03-L037,RP03-L038,RP03-L039
\item\label{refentry:REF049}R.M. Califf, D.A. Zarin, J.M. Kramer, R.E. Sherman, L.H. Aberle, and A. Tasneem, Characteristics of clinical trials registered in ClinicalTrials.gov, 2007--2010. J. Am. Med. Assoc. \textbf{307}(17), 1838--1847 (2012)\label{refend:REF049}
% REF050; 原书分章页 3; RP03-L040,RP03-L041
\item\label{refentry:REF050}A.C. Cameron, J.B. Gelbach, D.L. Miller, Boostrap-based improvements for inference with clustered errors. Rev. Econ. Stat. \textbf{90}(3), 414--427 (2008)\label{refend:REF050}
% REF051; 原书分章页 3; RP03-L042,RP03-L043
\item\label{refentry:REF051}M.J. Campbell, Cluster randomized trials, in \textit{Handbook} \textit{of} \textit{Epidemiology}, 2nd edn., ed. by W. Ahrens, I. Pigeot (Springer, Berlin, 2014)\label{refend:REF051}
% REF052; 原书分章页 3; RP03-L044,RP03-L045
\item\label{refentry:REF052}M.J. Campbell, A. Donner, N. Klar, Developments in cluster randomized trials and statistics in medicine. Stat. Med. \textbf{26}, 2--19 (2007)\label{refend:REF052}
% REF053; 原书分章页 4; RP04-L003,RP04-L004,RP04-L005
\item\label{refentry:REF053}M. Carone, I. Díaz, M.J. van der Laan, Higher-order targeted minimum loss-based estimation. Technical Report, Division of Biostatistics, University of California, Berkeley\label{refend:REF053}
% REF054; 原书分章页 4; RP04-L006,RP04-L007
\item\label{refentry:REF054}B. Chakraborty, E.E. Moodie, \textit{Statistical} \textit{Methods} \textit{for} \textit{Dynamic} \textit{Treatment} \textit{Regimes} (Springer, Berlin, Heidelberg, New York, 2013)\label{refend:REF054}
% REF055; 原书分章页 4; RP04-L008,RP04-L009
\item\label{refentry:REF055}B. Chakraborty, E.B. Laber, Y.-Q. Zhao, Inference about the expected performance of a data-driven dynamic treatment regime. Clin. Trials \textbf{11}(4), 408--417 (2014)\label{refend:REF055}
% REF056; 原书分章页 4; RP04-L010,RP04-L011
\item\label{refentry:REF056}A. Chambaz, tsml.cara.rct: targeted sequential minimum loss CARA RCT design and inference (2016). \url{https://github.com/achambaz/tsml.cara.rct}\label{refend:REF056}
% REF057; 原书分章页 4; RP04-L012,RP04-L013,RP04-L014
\item\label{refentry:REF057}A. Chambaz, P. Neuvial, Targeted, integrative search of associations between DNA copy number and gene expression, accounting for DNA methylation. Bioinformatics \textbf{31}(18), 3054--3056 (2015)\label{refend:REF057}
% REF058; 原书分章页 4; RP04-L015,RP04-L016,RP04-L017
\item\label{refentry:REF058}A. Chambaz, P. Neuvial, Targeted learning of a non-parametric variable importance measure of a continuous exposure (2016). \url{http://CRAN.R-project.org/package=tmle.npvi}\label{refend:REF058}
% REF059; 原书分章页 4; RP04-L018,RP04-L019
\item\label{refentry:REF059}A. Chambaz, M.J. van der Laan, Inference in targeted group-sequential covariate-adjusted randomized clinical trials. Scand. J. Stat. \textbf{41}(1), 104--140 (2014)\label{refend:REF059}
% REF060; 原书分章页 4; RP04-L020,RP04-L021,RP04-L022
\item\label{refentry:REF060}A. Chambaz, M.J. van der Laan, Targeting the optimal design in randomized clinical trials with binary outcomes and no covariate: theoretical study. Int. J. Biostat. \textbf{7}(1), Article 10 (2011a)\label{refend:REF060}
% REF061; 原书分章页 4; RP04-L023,RP04-L024,RP04-L025
\item\label{refentry:REF061}A. Chambaz, M.J. van der Laan, Targeting the optimal design in randomized clinical trials with binary outcomes and no covariate: simulation study. Int. J. Biostat. \textbf{7}(1), Article 11 (2011b)\label{refend:REF061}
% REF062; 原书分章页 4; RP04-L026,RP04-L027,RP04-L028,RP04-L029
\item\label{refentry:REF062}A. Chambaz, M.J. van der Laan, TMLE in adaptive group sequential covariate-adjusted RCTs, in \textit{Targeted} \textit{Learning:} \textit{Causal} \textit{Inference} \textit{for} \textit{Observational} \textit{and} \textit{Experimental} \textit{Data}, ed. by M.J. van der Laan, S. Rose (Springer, Berlin Heidelberg, New York, 2011c)\label{refend:REF062}
% REF063; 原书分章页 4; RP04-L030,RP04-L031,RP04-L032
\item\label{refentry:REF063}A. Chambaz, P. Neuvial, M.J. van der Laan, Estimation of a non-parametric variable importance measure of a continuous exposure. Electron. J. Stat. \textbf{6}, 1059--1099 (2012)\label{refend:REF063}
% REF064; 原书分章页 4; RP04-L033,RP04-L034,RP04-L035
\item\label{refentry:REF064}A. Chambaz, D. Choudat, C. Huber, J.C. Pairon, M.J. van der Laan, Analysis of the effect of occupational exposure to asbestos based on threshold regression modeling of case--control data. Biostatistics \textbf{15}(2), 327--340 (2014)\label{refend:REF064}
% REF065; 原书分章页 4; RP04-L036,RP04-L037,RP04-L038,RP04-L039
\item\label{refentry:REF065}A. Chambaz, M.J. van der Laan, W. Zheng, Targeted covariate-adjusted response-adaptive lasso-based randomized controlled trials, in \textit{Modern} \textit{Adaptive} \textit{Randomized} \textit{Clinical} \textit{Trials:} \textit{Statistical,} \textit{Operational,} \textit{and} \textit{Regulatory} \textit{Aspects}, ed. by A. Sverdlov (CRC Press, Boca Raton, 2015), pp. 345--368\label{refend:REF065}
% REF066; 原书分章页 4; RP04-L040,RP04-L041,RP04-L042
\item\label{refentry:REF066}A. Chambaz, W. Zheng, M.J. van der Laan, Targeted sequential design for targeted learning of the optimal treatment rule and its mean reward. Ann Stat. \textbf{45}(6), 1--28 (2017)\label{refend:REF066}
% REF067; 原书分章页 4; RP04-L043,RP04-L044,RP04-L045
\item\label{refentry:REF067}T. Chen, C. Guestrin, Xgboost: a scalable tree boosting system, in \textit{Proceedings} \textit{of} \textit{the} \textit{22nd} \textit{ACM} \textit{SIGKDD} \textit{International} \textit{Conference} \textit{on} \textit{Knowledge} \textit{Discovery} \textit{and} \textit{Data} \textit{Mining} (ACM, New York, 2016), pp. 785--794\label{refend:REF067}
% REF068; 原书分章页 5; RP05-L003,RP05-L004,RP05-L005
\item\label{refentry:REF068}O.Y. Chén, C. Crainiceanu, E.L. Ogburn, B.S. Caffo, T.D. Wager, M.A. Lindquist, High-dimensional multivariate mediation with application to neuroimaging data. Biostatistics (2017, in press)\label{refend:REF068}
% REF069; 原书分章页 5; RP05-L006,RP05-L007
\item\label{refentry:REF069}D.S. Choi, Estimation of monotone treatment effects in network experiments. ArXiv e-prints, August (2014)\label{refend:REF069}
% REF070; 原书分章页 5; RP05-L008,RP05-L009
\item\label{refentry:REF070}N.A. Christakis, J.H. Fowler, The spread of obesity in a large social network over 32 years. N. Engl. J. Med. \textbf{357}(4), 370--379 (2007)\label{refend:REF070}
% REF071; 原书分章页 5; RP05-L010,RP05-L011
\item\label{refentry:REF071}N.A. Christakis, J.H. Fowler, Social contagion theory: examining dynamic social networks and human behavior. Stat. Med. \textbf{32}(4), 556--577 (2013)\label{refend:REF071}
% REF072; 原书分章页 5; RP05-L012,RP05-L013
\item\label{refentry:REF072}W.G. Cochran, Analysis of covariance: its nature and uses. Biometrics \textbf{13}, 261--281 (1957)\label{refend:REF072}
% REF073; 原书分章页 5; RP05-L014,RP05-L015,RP05-L016
\item\label{refentry:REF073}E. Colantuoni, M. Rosenblum, Leveraging prognostic baseline variables to gain precision in randomized trials. Technical Report 263, Johns Hopkins University, Department of Biostatistics Working Papers (2015)\label{refend:REF073}
% REF074; 原书分章页 5; RP05-L017,RP05-L018
\item\label{refentry:REF074}S.R. Cole, E.A. Stuart, Generalizing evidence from randomized clinical trials to target populations: the ACTG 320 trial. Am. J. Epidemiol. \textbf{172}(1), 107--115 (2010)\label{refend:REF074}
% REF075; 原书分章页 5; RP05-L019,RP05-L020,RP05-L021,RP05-L022,RP05-L023
\item\label{refentry:REF075}S.R. Cole, M.A. Hernan, J.M. Robins, K. Anastos, J. Chmiel, R. Detels, C. Ervin, J. Feldman, R. Greenblatt, L. Kingsley, S. Lai, M. Young, M. Cohen, A. Munoz, Effect of highly active antiretroviral therapy on time to acquired immunodeficiency syndrome or death using marginal structural models. Am. J. Epidemiol. \textbf{158}(7), 687--694 (2003)\label{refend:REF075}
% REF076; 原书分章页 5; RP05-L024,RP05-L025
\item\label{refentry:REF076}D.R. Cox, P. McCullagh, Some aspects of analysis of covariance. Biometrics \textbf{38}(3), 541--561 (1982)\label{refend:REF076}
% REF077; 原书分章页 5; RP05-L026,RP05-L027
\item\label{refentry:REF077}K. Crammer, O. Dekel, J. Keshet, S. Shalev-Shwartz, Y. Singer, Online passive-aggressive algorithms. J. Mach. Learn. Res. \textbf{7}, 551--585 (2006)\label{refend:REF077}
% REF078; 原书分章页 5; RP05-L028,RP05-L029
\item\label{refentry:REF078}G.B. Dantzig, Discrete-variable extremum problems. Oper. Res. \textbf{5}(2), 266--288 (1957)\label{refend:REF078}
% REF079; 原书分章页 5; RP05-L030,RP05-L031,RP05-L032
\item\label{refentry:REF079}A.C. Davison, D.V. Hinkley, \textit{Bootstrap} \textit{methods} \textit{and} \textit{Their} \textit{Application}. Cambridge Series in Statistical and Probabilistic Mathematics, vol. 1 (Cambridge University Press, Cambridge, Cambridge, 1997)\label{refend:REF079}
% REF080; 原书分章页 5; RP05-L033,RP05-L034
\item\label{refentry:REF080}A.P. Dawid, V. Didelez, Identifying the consequences of dynamic treatment strategies: a decision-theoretic overview. Stat. Surv. \textbf{4}, 184--231 (2010)\label{refend:REF080}
% REF081; 原书分章页 5; RP05-L035,RP05-L036
\item\label{refentry:REF081}V.H. de la Peña, E. Giné, Decoupling, in \textit{Probability} \textit{and} \textit{its} \textit{Applications} (Springer, New York, 1999)\label{refend:REF081}
% REF082; 原书分章页 5; RP05-L037,RP05-L038
\item\label{refentry:REF082}L. Denby, C. Mallows, Variations on the histogram. J. Comput. Graph. Stat. \textbf{18}(1), 21--31 (2009)\label{refend:REF082}
% REF083; 原书分章页 5; RP05-L039,RP05-L040
\item\label{refentry:REF083}I. Díaz, M. van der Laan, Super learner-based conditional density estimation with application to marginal structural models. Int. J. Biostat. \textbf{7}(1), 38 (2011)\label{refend:REF083}
% REF084; 原书分章页 5; RP05-L041,RP05-L042
\item\label{refentry:REF084}I. Díaz, M. van der Laan, Population intervention causal effects based on stochastic interventions. Biometrics \textbf{68}(2), 541--549 (2012)\label{refend:REF084}
% REF085; 原书分章页 5; RP05-L043,RP05-L044
\item\label{refentry:REF085}I. Díaz, M.J. van der Laan, Assessing the causal effect of policies: an example using stochastic interventions. Int. J. Biostat. \textbf{9}(2), 161--174 (2013a)\label{refend:REF085}
% REF086; 原书分章页 5; RP05-L045,RP05-L046,RP05-L047
\item\label{refentry:REF086}I. Díaz, M.J. van der Laan, Sensitivity analysis for causal inference under unmeasured confounding and measurement error problems. Int. J. Biostat. \textbf{9}(2), 149--160 (2013b)\label{refend:REF086}
% REF087; 原书分章页 6; RP06-L003,RP06-L004
\item\label{refentry:REF087}I. Díaz, M. Carone, M.J. van der Laan, Second-order inference for the mean of a variable missing at random. Int. J. Biostat. \textbf{12}(1), 333--349 (2016)\label{refend:REF087}
% REF088; 原书分章页 6; RP06-L005,RP06-L006,RP06-L007
\item\label{refentry:REF088}I. Díaz, A. Hubbard, A. Decker, M. Cohen, Variable importance and prediction methods for longitudinal problems with missing variables. PLoS One \textbf{10}(3), e0120031 (2015)\label{refend:REF088}
% REF089; 原书分章页 6; RP06-L008,RP06-L009
\item\label{refentry:REF089}T.J. DiCiccio, J.P. Romano, A review of bootstrap confidence intervals. J. R. Stat. Soc. Ser. B (1988)\label{refend:REF089}
% REF090; 原书分章页 6; RP06-L010,RP06-L011,RP06-L012
\item\label{refentry:REF090}T.J. DiCiccio, J.P. Romano, Nonparametric confidence limits by resampling methods and least favorable families. Int. Stat. Rev./Revue Internationale de Statistique \textbf{58}(1), 59 (1990)\label{refend:REF090}
% REF091; 原书分章页 6; RP06-L013,RP06-L014,RP06-L015
\item\label{refentry:REF091}V. Didelez, A.P. Dawid, S. Geneletti, Direct and indirect effects of sequential treatments, in \textit{Proceedings} \textit{of} \textit{the} \textit{22nd} \textit{Annual} \textit{Conference} \textit{on} \textit{Uncertainty} \textit{in} \textit{Artificial} \textit{Intelligence} (2006), pp. 138--146\label{refend:REF091}
% REF092; 原书分章页 6; RP06-L016,RP06-L017
\item\label{refentry:REF092}P. Ding, T. VanderWeele, Sensitivity analysis without assumptions. Epidemiol. \textbf{27}(3), 368--377 (2016)\label{refend:REF092}
% REF093; 原书分章页 6; RP06-L018,RP06-L019
\item\label{refentry:REF093}A. Donner, N. Klar, \textit{Design} \textit{and} \textit{Analysis} \textit{of} \textit{Cluster} \textit{Randomization} \textit{Trials} \textit{in} \textit{Health} \textit{Research} (Arnold, London, 2000)\label{refend:REF093}
% REF094; 原书分章页 6; RP06-L020,RP06-L021
\item\label{refentry:REF094}J. Duchi, E. Hazan, Y. Singer, Adaptive subgradient methods for online learning and stochastic optimization. J. Mach. Learn. Res. \textbf{12}, 2121--2159 (2011)\label{refend:REF094}
% REF095; 原书分章页 6; RP06-L022,RP06-L023,RP06-L024,RP06-L025,RP06-L026
\item\label{refentry:REF095}W. Duckworth, C. Abraira, T. Moritz, D. Reda, N. Emanuele, P.D. Reaven, F.J. Zieve, J. Marks, S.N. Davis, R. Hayward, S.R. Warren, S. Goldman, M. McCarren, M.E. Vitek, W.G. Henderson, G.D. Huang for the VADT Investigators, Glucose control and vascular complications in veterans with type 2 diabetes. N. Engl. J. Med. \textbf{360}, 129--39 (2009a)\label{refend:REF095}
% REF096; 原书分章页 6; RP06-L027,RP06-L028
\item\label{refentry:REF096}W. Duckworth et al., Glucose control and vascular complications in veterans with type 2 diabetes. N. Engl. J. Med. \textbf{360}(2), 129--139 (2009b)\label{refend:REF096}
% REF097; 原书分章页 6; RP06-L029,RP06-L030,RP06-L031
\item\label{refentry:REF097}S. Dudoit, M.J. van der Laan, Asymptotics of cross-validated risk estimation in estimator selection and performance assessment. Stat. Methodol. \textbf{2}(2), 131--154 (2005)\label{refend:REF097}
% REF098; 原书分章页 6; RP06-L032,RP06-L033
\item\label{refentry:REF098}F. Eberhardt, R. Scheines, Interventions and causal inference. Department of Philosophy. Paper 415 (2006)\label{refend:REF098}
% REF099; 原书分章页 6; RP06-L034,RP06-L035
\item\label{refentry:REF099}B. Efron, Better bootstrap confidence intervals. J. Am. Stat. Assoc. \textbf{82}(397), 171--185 (1987)\label{refend:REF099}
% REF100; 原书分章页 6; RP06-L036,RP06-L037
\item\label{refentry:REF100}B. Efron, R.J. Tibshirani, \textit{An} \textit{Introduction} \textit{to} \textit{the} \textit{Bootstrap} (Chapman \& Hall, Boca Raton, 1993)\label{refend:REF100}
% REF101; 原书分章页 6; RP06-L038,RP06-L039
\item\label{refentry:REF101}U. Einmahl, D.M. Mason, An empirical process approach to the uniform consistency of kernel-type function estimators. J. Theor. Probab. \textbf{13}(1) 1--37 (2000)\label{refend:REF101}
% REF102; 原书分章页 6; RP06-L040,RP06-L041
\item\label{refentry:REF102}U. Einmahl, D.M. Mason, Uniform in bandwidth consistency of kernel-type function estimators. Ann. Stat. \textbf{33}(3), 1380--1403 (2005)\label{refend:REF102}
% REF103; 原书分章页 6; RP06-L042,RP06-L043
\item\label{refentry:REF103}European Medicines Agency, Guideline on adjustment for baseline covariates in clinical trials. London, February (2015)\label{refend:REF103}
% REF104; 原书分章页 6; RP06-L044,RP06-L045
\item\label{refentry:REF104}J.P. Fine, R.J. Gray, A proportional hazards model for the subdistribution of a competing risk. J. Am. Stat. Assoc. \textbf{94}(446), 496--509 (1999)\label{refend:REF104}
% REF105; 原书分章页 6; RP06-L046,RP06-L047
\item\label{refentry:REF105}M. Finster, M. Wood, The Apgar score has survived the test of time. Anesthesiology \textbf{102}(4), 855--857 (2005)\label{refend:REF105}
% REF106; 原书分章页 7; RP07-L003,RP07-L004
\item\label{refentry:REF106}R.A. Fisher, \textit{Statistical} \textit{Methods} \textit{for} \textit{Research} \textit{Workers}, 4th edn. (Oliver and Boyd Ltd., Edinburgh, 1932)\label{refend:REF106}
% REF107; 原书分章页 7; RP07-L005
\item\label{refentry:REF107}R.A. Fisher, \textit{The} \textit{Design} \textit{of} \textit{Experiments}, (Oliver and Boyd Ltd, London, 1935)\label{refend:REF107}
% REF108; 原书分章页 7; RP07-L006,RP07-L007,RP07-L008
\item\label{refentry:REF108}C.E. Frangakis, T. Qian, Z. Wu, I. Diaz, Deductive derivation and Turing-computerization of semiparametric efficient estimation. Biometrics \textbf{71}(4), 867--874 (2015)\label{refend:REF108}
% REF109; 原书分章页 7; RP07-L009,RP07-L010,RP07-L011
\item\label{refentry:REF109}L.S. Freedman, M.H. Gail, S.B. Green, D.K. Corle, The COMMIT Research Group, The Efficiency of the matched-pairs design of the community intervention trial for smoking cessation (COMMIT). Control. Clin. Trials \textbf{18}(2), 131--139 (1997)\label{refend:REF109}
% REF110; 原书分章页 7; RP07-L012,RP07-L013
\item\label{refentry:REF110}J.H. Friedman, Multivariate adaptive regression splines. Ann. Stat. \textbf{19}(1), 1--141 (1991)\label{refend:REF110}
% REF111; 原书分章页 7; RP07-L014,RP07-L015
\item\label{refentry:REF111}J.H. Friedman, Greedy function approximation: a gradient boosting machine. Ann. Stat. \textbf{29}, 1189--1232 (2001)\label{refend:REF111}
% REF112; 原书分章页 7; RP07-L016,RP07-L017
\item\label{refentry:REF112}J.H. Friedman, T.J. Hastie, R.J. Tibshirani, Glmnet: lasso and elastic-net regularized generalized linear models (2010). \url{http://CRAN.R-project.org/package=glmnet}\label{refend:REF112}
% REF113; 原书分章页 7; RP07-L018,RP07-L019
\item\label{refentry:REF113}K.J. Friston, L. Harrison, W. Penny, Dynamic causal modelling. Neuroimage \textbf{19}(4), 1273--1302 (2003)\label{refend:REF113}
% REF114; 原书分章页 7; RP07-L020,RP07-L021
\item\label{refentry:REF114}K. Friston, R. Moran, A.K. Seth, Analysing connectivity with granger causality and dynamic causal modelling. Curr. Opin. Neurobiol. \textbf{23}(2), 172--178 (2013)\label{refend:REF114}
% REF115; 原书分章页 7; RP07-L022,RP07-L023
\item\label{refentry:REF115}W.J. Fu, Penalized regressions: the bridge versus the lasso. J. Comput. Graph. Stat. \textbf{7}(3), 397--416 (1998)\label{refend:REF115}
% REF116; 原书分章页 7; RP07-L024
\item\label{refentry:REF116}P. Galison, \textit{How} \textit{Experiments} \textit{End} (University of Chicago Press, Chicago, 1987)\label{refend:REF116}
% REF117; 原书分章页 7; RP07-L025,RP07-L026,RP07-L027,RP07-L028
\item\label{refentry:REF117}J.J. Gaynor, E.J. Feuer, C.C. Tan, D.H. Wu, C.R. Little, D.J. Straus, B.D. Clarkson, M.F. Brennan, On the use of cause-specific failure and conditional failure probabilities: examples from clinical oncology data. J. Am. Stat. Assoc. \textbf{88}(422), 400--409 (1993)\label{refend:REF117}
% REF118; 原书分章页 7; RP07-L029,RP07-L030
\item\label{refentry:REF118}A. Gelman, C. Shalizi, Philosophy and the practice of bayesian statistics. Br. J. Math. Stat. Psychol. 66(1), 8--38 (2013)\label{refend:REF118}
% REF119; 原书分章页 7; RP07-L031,RP07-L032,RP07-L033
\item\label{refentry:REF119}A. Gelman, Y.-S. Su, M. Yajima, J. Hill, M.G. Pittau, J. Kerman, T. Zheng, Arm: data analysis using regression and multilevel/hierarchical models (2010). \url{http://CRAN.R-project.org/package=arm}\label{refend:REF119}
% REF120; 原书分章页 7; RP07-L034,RP07-L035
\item\label{refentry:REF120}H.C. Gerstein et al., Effects of intensive glucose lowering in type 2 diabetes. N. Engl. J. Med. \textbf{358}(24), 2545--2559 (2008)\label{refend:REF120}
% REF121; 原书分章页 7; RP07-L036,RP07-L037
\item\label{refentry:REF121}G. Gigerenzer, \textit{The} \textit{Empire} \textit{of} \textit{Chance:} \textit{How} \textit{Probability} \textit{Changed} \textit{Science} \textit{and} \textit{Everyday} \textit{Life} (Cambridge University Press, Cambridge, 1989)\label{refend:REF121}
% REF122; 原书分章页 7; RP07-L038,RP07-L039,RP07-L040
\item\label{refentry:REF122}P.B. Gilbert, Comparison of competing risks failure time methods and time-independent methods for assessing strain variations in vaccine protection. Stat. Med. \textbf{19}(22), 3065--3086 (2000)\label{refend:REF122}
% REF123; 原书分章页 7; RP07-L041,RP07-L042,RP07-L043
\item\label{refentry:REF123}P.B. Gilbert, S.G. Self, M.A. Ashby, Statistical methods for assessing differential vaccine protection against human immunodeficiency virus types. Biometrics \textbf{54}(3), 799--814 (1998)\label{refend:REF123}
% REF124; 原书分章页 7; RP07-L044,RP07-L045,RP07-L046
\item\label{refentry:REF124}P.B. Gilbert, S.G. Self, M. Rao, A. Naficy, J. Clemens, Sieve analysis: methods for assessing from vaccine trial data how vaccine efficacy varies with genotypic and phenotypic pathogen variation. J. Clin. Epidemiol. \textbf{54}(1), 68--85 (2001)\label{refend:REF124}
% REF125; 原书分章页 8; RP08-L003,RP08-L004
\item\label{refentry:REF125}R.D. Gill, Non- and semiparametric maximum likelihood estimators and the von Mises method (Part 1). Scand. J. Stat. \textbf{16}, 91--128 (1989)\label{refend:REF125}
% REF126; 原书分章页 8; RP08-L005,RP08-L006
\item\label{refentry:REF126}R.D. Gill, J.M. Robins, Causal inference in complex longitudinal studies: continuous case. Ann. Stat. \textbf{29}(6), 1785--1811 (2001)\label{refend:REF126}
% REF127; 原书分章页 8; RP08-L007,RP08-L008,RP08-L009
\item\label{refentry:REF127}R.D. Gill, M.J. van der Laan, J.A. Wellner, Inefficient estimators of the bivariate survival function for three models. Ann. l'Institut Henri Poincaré \textbf{31}(3), 545--597 (1995)\label{refend:REF127}
% REF128; 原书分章页 8; RP08-L010,RP08-L011,RP08-L012
\item\label{refentry:REF128}Y. Goldberg, R. Song, D. Zeng, M.R. Kosorok, Comment on ``Dynamic treatment regimes: technical challenges and applications''. Electron. J. Stat. \textbf{8}, 1290--1300 (2014)\label{refend:REF128}
% REF129; 原书分章页 8; RP08-L013,RP08-L014,RP08-L015
\item\label{refentry:REF129}N. Grambauer, M. Schumacher, J. Beyersmann, Proportional subdistribution hazards modeling offers a summary analysis, even if misspecified. Stat. Med. \textbf{29}(7--8), 875--884 (2010)\label{refend:REF129}
% REF130; 原书分章页 8; RP08-L016,RP08-L017
\item\label{refentry:REF130}R. Greevy, B. Lu, J.H. Silber, P. Rosenbaum, Optimal multivariate matching before randomization. Biostatistics \textbf{5}(2), 263--275 (2004)\label{refend:REF130}
% REF131; 原书分章页 8; RP08-L018,RP08-L019
\item\label{refentry:REF131}U. Grömping, Variable importance assessment in regression: linear regression versus random forest. Am. Stat. \textbf{63}(4) (2009)\label{refend:REF131}
% REF132; 原书分章页 8; RP08-L020,RP08-L021,RP08-L022,RP08-L023,RP08-L024
\item\label{refentry:REF132}H. Grosskurth, F. Mosha, J. Todd, E. Mwijarubi, A. Klokke, K. Senkoro, P. Mayaud, J. Changalucha, A. Nicoll, G. ka-Gina, J. Newell, K. Mugeye, D. Mabey, R. Hayes, Impact of improved treatment of sexually transmitted diseases on HIV infection in rural Tanzania: randomised controlled trial. Lancet \textbf{346}(8974), 530--536 (1995)\label{refend:REF132}
% REF133; 原书分章页 8; RP08-L025,RP08-L026,RP08-L027
\item\label{refentry:REF133}S. Gruber, M.J. van der Laan, An application of collaborative targeted maximum likelihood estimation in causal inference and genomics. Int. J. Biostat. \textbf{6}(1) (2010a)\label{refend:REF133}
% REF134; 原书分章页 8; RP08-L028,RP08-L029
\item\label{refentry:REF134}S. Gruber, M.J. van der Laan, A targeted maximum likelihood estimator of a causal effect on a bounded continuous outcome. Int. J. Biostat. \textbf{6}(1), Article 26 (2010b)\label{refend:REF134}
% REF135; 原书分章页 8; RP08-L030,RP08-L031
\item\label{refentry:REF135}S. Gruber, M.J. van der Laan, tmle: an R package for targeted maximum likelihood estimation. J. Stat. Softw. 51(13) (2012a)\label{refend:REF135}
% REF136; 原书分章页 8; RP08-L032,RP08-L033
\item\label{refentry:REF136}S. Gruber, M.J. van der Laan, Targeted minimum loss based estimator that outperforms a given estimator. Int. J. Biostat. \textbf{8}(1), (2012b)\label{refend:REF136}
% REF137; 原书分章页 8; RP08-L034,RP08-L035
\item\label{refentry:REF137}I. Hacking, \textit{The} \textit{Emergence} \textit{of} \textit{Probability} (Cambridge University Press, Cambridge, 1975)\label{refend:REF137}
% REF138; 原书分章页 8; RP08-L036,RP08-L037
\item\label{refentry:REF138}I. Hacking, \textit{The} \textit{Taming} \textit{of} \textit{Chance} \textit{(1990)} (Cambridge University Press, Cambridge, 1990)\label{refend:REF138}
% REF139; 原书分章页 8; RP08-L038,RP08-L039
\item\label{refentry:REF139}D.M. Hafeman, T.J. VanderWeele, Alternative assumptions for the identification of direct and indirect effects. Epidemiology \textbf{22}, 753--764 (2010)\label{refend:REF139}
% REF140; 原书分章页 8; RP08-L040,RP08-L041
\item\label{refentry:REF140}J. Hahn, On the role of the propensity score in efficient semiparametric estimation of average treatment effects. Econometrica \textbf{2}, 315--331 (1998)\label{refend:REF140}
% REF141; 原书分章页 8; RP08-L042,RP08-L043
\item\label{refentry:REF141}J. Hajek, Asymptotic theory of rejective sampling with varying probabilities from a finite population. Ann. Math. Stat. \textbf{35}(4), 1491--1523, 12 (1964)\label{refend:REF141}
% REF142; 原书分章页 8; RP08-L044,RP08-L045
\item\label{refentry:REF142}P Hall, Theoretical comparison of bootstrap confidence intervals. Ann. Stat. \textbf{16}, 927--953 (1988)\label{refend:REF142}
% REF143; 原书分章页 8; RP08-L046,RP08-L047
\item\label{refentry:REF143}P. Hall, \textit{The} \textit{Bootstrap} \textit{and} \textit{Edgeworth} \textit{Expansion}. Springer Series in Statistics (Springer, New York, NY, 1992)\label{refend:REF143}
% REF144; 原书分章页 9; RP09-L003,RP09-L004
\item\label{refentry:REF144}M.E. Halloran, C.J. Struchiner, Causal inference in infectious diseases. Epidemiology \textbf{6}(2), 142--151 (1995)\label{refend:REF144}
% REF145; 原书分章页 9; RP09-L005,RP09-L006,RP09-L007,RP09-L008
\item\label{refentry:REF145}S.M. Hammer, M.E. Sobieszczyk, H. Janes, S.T. Karuna, M.J. Mulligan, D. Grove, B.A. Koblin, S.P. Buchbinder, M.C. Keefer, G.D. Tomaras, Efficacy trial of a DNA/rAd5 HIV-1 preventive vaccine. N. Engl. J. Med. \textbf{369}(22), 2083--2092 (2013)\label{refend:REF145}
% REF146; 原书分章页 9; RP09-L009,RP09-L010
\item\label{refentry:REF146}S. Haneuse, A. Rotnitzky, Estimation of the effect of interventions that modify the received treatment. Stat. Med. (2013)\label{refend:REF146}
% REF147; 原书分章页 9; RP09-L011,RP09-L012,RP09-L013
\item\label{refentry:REF147}M. Hanif, K.R.W. Brewer, Sampling with unequal probabilities without replacement: a review. International Statistical Review/Revue Internationale de Statistique, pp. 317--335 (1980)\label{refend:REF147}
% REF148; 原书分章页 9; RP09-L014,RP09-L015,RP09-L016,RP09-L017
\item\label{refentry:REF148}E. Hartman, R. Grieve, R. Ramsahai, J.S. Sekhon, From sample average treatment effect to population average treatment effect on the treated: combining experimental with observational studies to estimate population treatment effects. J. R. Stat. Soc. Ser. A \textbf{178}(3), 757--778 (2015)\label{refend:REF148}
% REF149; 原书分章页 9; RP09-L018,RP09-L019
\item\label{refentry:REF149}T. Hastie, gam: generalized additive models (2011) \url{http://CRAN.R-project.org/package=gam}\label{refend:REF149}
% REF150; 原书分章页 9; RP09-L020,RP09-L021,RP09-L022
\item\label{refentry:REF150}T.J. Hastie, R.J. Tibshirani, J.H. Friedman, \textit{The} \textit{Elements} \textit{of} \textit{Statistical} \textit{Learning:} \textit{Data} \textit{Mining,} \textit{Inference,} \textit{and} \textit{Prediction} (Springer, Berlin Heidelberg New York, 2001)\label{refend:REF150}
% REF151; 原书分章页 9; RP09-L023,RP09-L024
\item\label{refentry:REF151}R.J. Hayes, L.H. Moulton, \textit{Cluster} \textit{Randomised} \textit{Trials}. (Chapman \& Hall/CRC, Boca Raton, 2009)\label{refend:REF151}
% REF152; 原书分章页 9; RP09-L025,RP09-L026
\item\label{refentry:REF152}M.A. Hearst, S.T Dumais, E. Osman, J. Platt, B. Scholkopf. Support vector machines. IEEE Intell. Syst. Appl. \textbf{13}(4), 18--28 (1998)\label{refend:REF152}
% REF153; 原书分章页 9; RP09-L027,RP09-L028,RP09-L029
\item\label{refentry:REF153}M.A. Hernan, B.A. Brumback, J.M. Robins, Estimating the causal effect of zidovudine on CD4 count with a marginal structural model for repeated measures. Stat. Med. \textbf{21}, 1689--1709 (2002)\label{refend:REF153}
% REF154; 原书分章页 9; RP09-L030,RP09-L031,RP09-L032
\item\label{refentry:REF154}M.A. Hernan, B. Brumback, J.M. Robins, Marginal structural models to estimate the causal effect of zidovudine on the survival of HIV-positive men. Epidemiology \textbf{11}(5), 561--570 (2000)\label{refend:REF154}
% REF155; 原书分章页 9; RP09-L033,RP09-L034,RP09-L035
\item\label{refentry:REF155}M.A. Hernan, E. Lanoy, D. Costagliola, J.M. Robins, Comparison of dynamic treatment regimes via inverse probability weighting. Basic Clin. Pharmacol. \textbf{98}, 237--242 (2006)\label{refend:REF155}
% REF156; 原书分章页 9; RP09-L036,RP09-L037
\item\label{refentry:REF156}R. Holiday, What the failed \$1m Netflix prize says about business advice. Forbes (2012)\label{refend:REF156}
% REF157; 原书分章页 9; RP09-L038,RP09-L039,RP09-L040
\item\label{refentry:REF157}R.R. Holman, S.K. Paul, M.A. Bethel, D.R. Matthews, H.A. Neil, 10-year follow-up of intensive glucose control in type 2 diabetes. N. Engl. J. Med. \textbf{359}, 1577--89 (2008)\label{refend:REF157}
% REF158; 原书分章页 9; RP09-L041,RP09-L042,RP09-L043
\item\label{refentry:REF158}J.L. Horowitz, C.F. Manski, Nonparametric analysis of randomized experiments with missing covariate and outcome data. J. Am. Stat. Assoc. \textbf{95}(449), 77--84 (2000)\label{refend:REF158}
% REF159; 原书分章页 9; RP09-L044,RP09-L045
\item\label{refentry:REF159}D.G. Horvitz, D.J. Thompson, A generalization of sampling without replacement from a finite universe. J. Am. Stat. Assoc. \textbf{47}, 663--685 (1952)\label{refend:REF159}
% REF160; 原书分章页 9; RP09-L046,RP09-L047
\item\label{refentry:REF160}J.B. Hough, M. Krishnapur, Y. Peres, B. Virág, Determinantal processes and independence. Probab. Surv. \textbf{3}, 206--229 (2006)\label{refend:REF160}
% REF161; 原书分章页 10; RP10-L003,RP10-L004
\item\label{refentry:REF161}F. Hu, W.F. Rosenberger, \textit{The} \textit{Theory} \textit{of} \textit{Response} \textit{Adaptive} \textit{Randomization} \textit{in} \textit{Clinical} \textit{Trials} (Wiley, New York, 2006)\label{refend:REF161}
% REF162; 原书分章页 10; RP10-L005,RP10-L006,RP10-L007,RP10-L008
\item\label{refentry:REF162}A.E. Hubbard, M.J. van der Laan, Mining with inference: data adaptive target parameters, in \textit{Handbook} \textit{of} \textit{Big} \textit{Data}. Chapman-Handbooks-Statistical-Methods, ed. by P. Buhlmann, P. Drineas, M. Kane, M.J. van der Laan (Chapman \& Hall/CRC, Boca Raton, 2016)\label{refend:REF162}
% REF163; 原书分章页 10; RP10-L009,RP10-L010,RP10-L011,RP10-L012,RP10-L013
\item\label{refentry:REF163}A.E. Hubbard, I Diaz Munoz, A. Decker, J.B. Holcomb, M.A. Schreiber, E.M. Bulger, K.J. Brasel, E.E. Fox, D.J. del Junco, C.E. Wade et al., Time-dependent prediction and evaluation of variable importance using superlearning in high-dimensional clinical data. J. Trauma-Injury Infect. Crit. Care \textbf{75}(1), S53--S60 (2013)\label{refend:REF163}
% REF164; 原书分章页 10; RP10-L014,RP10-L015
\item\label{refentry:REF164}A.E. Hubbard, S. Kherad-Pajouh, M.J. van der Laan, Statistical inference for data adaptive target parameters. Int. J. Biostat. \textbf{12}(1), 3--19 (2016)\label{refend:REF164}
% REF165; 原书分章页 10; RP10-L016,RP10-L017
\item\label{refentry:REF165}M.G. Hudgens, M.E. Halloran, Toward causal inference with interference. J. Am. Stat. Assoc. \textbf{103}(482), 832--842 (2008)\label{refend:REF165}
% REF166; 原书分章页 10; RP10-L018
\item\label{refentry:REF166}I.A. Ibragimov, R.Z. Khasminskii, \textit{Statistical} \textit{Estimation} (Springer, Berlin, 1981)\label{refend:REF166}
% REF167; 原书分章页 10; RP10-L019,RP10-L020
\item\label{refentry:REF167}ICH Harmonised Tripartite Guideline, Statistical principles for clinical trials E9, February (1998)\label{refend:REF167}
% REF168; 原书分章页 10; RP10-L021,RP10-L022
\item\label{refentry:REF168}K. Imai, Variance identification and efficiency analysis in randomized experiments under the matched-pair design. Stat. Med. \textbf{27}(24), 4857--4873 (2008)\label{refend:REF168}
% REF169; 原书分章页 10; RP10-L023,RP10-L024,RP10-L025
\item\label{refentry:REF169}K. Imai, G. King, C. Nall, The essential role of pair matching in cluster-randomized experiments, with application to the Mexican universal health insurance evaluation. Stat. Sci. \textbf{24}(1), 29--53 (2009)\label{refend:REF169}
% REF170; 原书分章页 10; RP10-L026,RP10-L027
\item\label{refentry:REF170}K. Imai, L. Keele, D. Tingley, A general approach to causal mediation analysis. Psychol methods \textbf{15}(4), 309--334 (2010a)\label{refend:REF170}
% REF171; 原书分章页 10; RP10-L028,RP10-L029
\item\label{refentry:REF171}K. Imai, L. Keele, T. Yamamoto, Identification, inference and sensitivity analysis for causal mediation effects. Stat. Sci. \textbf{25}(1), 51--71 (2010b)\label{refend:REF171}
% REF172; 原书分章页 10; RP10-L030,RP10-L031
\item\label{refentry:REF172}G.W. Imbens, Nonparametric estimation of average treatment effects under exogeneity: a review. Rev. Econ. Stat. \textbf{86}(1), 4--29 (2004)\label{refend:REF172}
% REF173; 原书分章页 10; RP10-L032,RP10-L033
\item\label{refentry:REF173}G.W. Imbens, Experimental design for unit and cluster randomized trials. Technical Report. NBER Working Paper (2011)\label{refend:REF173}
% REF174; 原书分章页 10; RP10-L034,RP10-L035
\item\label{refentry:REF174}G. Imbens, D.B. Rubin, \textit{Causal} \textit{Inference} \textit{for} \textit{Statistics,} \textit{Social,} \textit{and} \textit{Biomedical} \textit{Sciences} (Cambridge University Press, New York, 2015)\label{refend:REF174}
% REF175; 原书分章页 10; RP10-L036,RP10-L037
\item\label{refentry:REF175}J.P. Ioannidis, Why most discovered true associations are inflated. Epidemiology \textbf{19}(5), 640--648 (2008)\label{refend:REF175}
% REF176; 原书分章页 10; RP10-L038,RP10-L039,RP10-L040,RP10-L041,RP10-L042,RP10-L043,RP10-L044
\item\label{refentry:REF176}F. Ismail-Beigi, T. Craven, M.A. Banerji, J. Basile, J. Calles, R.M. Cohen, R. Cuddihy, W.C Cushman, S. Genuth, R.H. Grimm, B.P. Hamilton, B. Hoogwerf, D. Karl, L. Katz, A. Krikorian, P. O'Connor, R. Pop-Busui, U. Schubart, D. Simmons, H. Taylor, A. Thomas, D. Weiss, I. Hramiak for the ACCORD trial group, Effect of intensive treatment of hyperglycaemia on microvascular outcomes in type 2 diabetes: an analysis of the ACCORD randomised trial. Lancet \textbf{376}, 419--430 (2010)\label{refend:REF176}
% REF177; 原书分章页 10; RP10-L045,RP10-L046,RP10-L047
\item\label{refentry:REF177}Joint National Committee, The fifth report of the joint national committee on detection, evaluation, and treatment of high blood pressure (JNC V). Arch. Intern. Med. \textbf{153}(2), 154--183 (1993)\label{refend:REF177}
% REF178; 原书分章页 11; RP11-L003,RP11-L004,RP11-L005
\item\label{refentry:REF178}B.C. Kahn, V. Jairath, C.J. Doré, T.P. Morris, The risks and rewards of covariate adjustment in randomized trials: an assessment of 12 outcomes from 8 studies. Trials \textbf{15}(139), 1--7 (2014)\label{refend:REF178}
% REF179; 原书分章页 11; RP11-L006,RP11-L007
\item\label{refentry:REF179}R.M. Karp, \textit{Reducibility} \textit{Among} \textit{Combinatorial} \textit{Problems} (Springer, New York, Berlin, Heidelberg, 1972)\label{refend:REF179}
% REF180; 原书分章页 11; RP11-L008,RP11-L009,RP11-L010
\item\label{refentry:REF180}S. Keleş, M.J. van der Laan, S. Dudoit, Asymptotically optimal model selection method for regression on censored outcomes. Technical Report, Division of Biostatistics, University of California, Berkeley (2002)\label{refend:REF180}
% REF181; 原书分章页 11; RP11-L011,RP11-L012
\item\label{refentry:REF181}E.H. Kennedy, Z. Ma, M.D. McHugh, D.S. Small, Nonparametric methods for doubly robust estimation of continuous treatment effects. ArXiv e-prints (2015)\label{refend:REF181}
% REF182; 原书分章页 11; RP11-L013,RP11-L014,RP11-L015
\item\label{refentry:REF182}R. Kessler, S. Rose, K. Koenen et al., How well can post-traumatic stress disorder be predicted from pre-trauma risk factors? an exploratory study in the who world mental health surveys. World Psychiatry \textbf{13}(3), 265--274 (2014)\label{refend:REF182}
% REF183; 原书分章页 11; RP11-L016,RP11-L017
\item\label{refentry:REF183}D. Kibler, D.W. Aha, M.K. Albert, Instance-based prediction of real-valued attributes. Comput. Intell. \textbf{5}, 51 (1989)\label{refend:REF183}
% REF184; 原书分章页 11; RP11-L018,RP11-L019
\item\label{refentry:REF184}J. Kivinen, A.J. Smola, R.C. Williamson, Online learning with kernels. IEEE Trans. Signal Process. \textbf{52}(8), 2165--2176 (2004)\label{refend:REF184}
% REF185; 原书分章页 11; RP11-L020,RP11-L021
\item\label{refentry:REF185}N. Klar, A. Donner, The merits of matching in community intervention trials: a cautionary tale. Stat. Med. \textbf{16}(15), 1753--1764 (1997)\label{refend:REF185}
% REF186; 原书分章页 11; RP11-L022,RP11-L023
\item\label{refentry:REF186}D.C. Knill, A. Pouget, The Bayesian brain: the role of uncertainty in neural coding and computation. Trends Neurosci. \textbf{27}(12), 712--719 (2004)\label{refend:REF186}
% REF187; 原书分章页 11; RP11-L024,RP11-L025,RP11-L026,RP11-L027
\item\label{refentry:REF187}K. Korb, L. Hope, A. Nicholson, K. Axnick, Varieties of causal intervention. in \textit{PRICAI} \textit{2004:} \textit{Trends} \textit{in} \textit{Artificial} \textit{Intelligence}, ed. by C. Zhang, H.W. Guesgen, W.-K. Yeap. Lecture Notes in Computer Science, vol. 3157 (Springer, Berlin, Heidelberg, 2004), pp. 322--331\label{refend:REF187}
% REF188; 原书分章页 11; RP11-L028,RP11-L029
\item\label{refentry:REF188}B. Korte, J. Vygen, \textit{Combinatorial} \textit{Optimization}, 5th edn. (Springer, Berlin, Heidelberg, New York, 2012)\label{refend:REF188}
% REF189; 原书分章页 11; RP11-L030,RP11-L031,RP11-L032,RP11-L033,RP11-L034
\item\label{refentry:REF189}M.S. Kramer, B. Chalmers, E.D. Hodnett, Z. Sevkovskaya, I. Dzikovich, S. Shapiro, J.P. Collet, I. Vanilovich, I. Mezen, T. Ducruet, G. Shishko, V. Zubovich, D. Mknuik, E. Gluchanina, V. Dombrovskiy, A. Ustinovitch, T. Kot, N. Bogdanovich, L. Ovchinikova, E. Helsing, PROmotion of breastfeeding intervention trial (PROBIT). J. Am. Med. Assoc. \textbf{285}(4), 413--420 (2001)\label{refend:REF189}
% REF190; 原书分章页 11; RP11-L035,RP11-L036,RP11-L037
\item\label{refentry:REF190}M.S. Kramer, T. Guo, R.W. Platt, S. Shapiro, J.P. Collet, B. Chalmers, E. Hodnett, Z. Sevkovskaya, I. Dzikovich, I. Vanilovich, Breastfeeding and infant growth: biology or bias? Pediatrics \textbf{110}(2), 343--347 (2002)\label{refend:REF190}
% REF191; 原书分章页 11; RP11-L038,RP11-L039
\item\label{refentry:REF191}L. Krüger, L. Daston, M. Heidelberger, G. Gigerenzer, M.S. Morgan, \textit{The} \textit{Probabilistic} \textit{Revolution}. (MIT Press, Cambridge, 1987)\label{refend:REF191}
% REF192; 原书分章页 11; RP11-L040,RP11-L041,RP11-L042
\item\label{refentry:REF192}L. Kunz, S. Rose, D. Spiegelman, S.-L. Normand, Causal inference methods in comparative effectiveness research, in \textit{Methods} \textit{in} \textit{Comparative} \textit{Effectiveness} \textit{Research}, ed. by C. Gatsonis, S.C. Morton (Chapman \& Hall, Boca Raton, 2017)\label{refend:REF192}
% REF193; 原书分章页 11; RP11-L043,RP11-L044,RP11-L045
\item\label{refentry:REF193}E.B. Laber, D.J. Lizotte, M. Qian, W.E. Pelham, S.A. Murphy, Dynamic treatment regimes: Technical challenges and applications. Electron. J. Stat. \textbf{8}(1), 1225--1272 (2014a)\label{refend:REF193}
% REF194; 原书分章页 12; RP12-L003,RP12-L004,RP12-L005
\item\label{refentry:REF194}E.B. Laber, D.J. Lizotte, M. Qian, W.E. Pelham, S.A. Murphy, Rejoinder of ``Dynamic treatment regimes: technical challenges and applications''. Electron. J. Stat. \textbf{8}(1), 1312--1321 (2014b)\label{refend:REF194}
% REF195; 原书分章页 12; RP12-L006,RP12-L007
\item\label{refentry:REF195}J. Langford, L. Li, T. Zhang, Sparse online learning via truncated gradient. J. Mach. Learn. Res. \textbf{10}, 777--801 (2009)\label{refend:REF195}
% REF196; 原书分章页 12; RP12-L008,RP12-L009
\item\label{refentry:REF196}P. Lavori, R. Dawson, Adaptive treatment strategies in chronic disease. Annu. Rev. Med. \textbf{59}, 443--453 (2008)\label{refend:REF196}
% REF197; 原书分章页 12; RP12-L010,RP12-L011
\item\label{refentry:REF197}P.W. Lavori, R. Dawson, A design for testing clinical strategies: Biased adaptive within-subject randomization. J. R. Stat. Soc. Ser. A \textbf{163} 29--38 (2000)\label{refend:REF197}
% REF198; 原书分章页 12; RP12-L012,RP12-L013
\item\label{refentry:REF198}D. Lazer, R. Kennedy, What we can learn from the epic failure of Google flu trends. Wired (2015)\label{refend:REF198}
% REF199; 原书分章页 12; RP12-L014,RP12-L015,RP12-L016
\item\label{refentry:REF199}S.D. Lendle, M.J. van der Laan, Identification and efficient estimation of the natural direct effect among the untreated. Technical Report, Division of Biostatistics, University of California, Berkeley (2011)\label{refend:REF199}
% REF200; 原书分章页 12; RP12-L017,RP12-L018,RP12-L019
\item\label{refentry:REF200}S.D. Lendle, B. Fireman, M.J. van der Laan, Balancing score adjusted targeted minimum loss-based estimation. Technical Report, Division of Biostatistics, University of California, Berkeley (2013)\label{refend:REF200}
% REF201; 原书分章页 12; RP12-L020,RP12-L021,RP12-L022
\item\label{refentry:REF201}S. Lendle, J. Schwab, M.L. Petersen, M.J. van der Laan, ltmle: an R package for implementing targeted minimum loss-based estimation for longitudinal data. J. Stat. Softw. \textbf{81}(1) (2017)\label{refend:REF201}
% REF202; 原书分章页 12; RP12-L023,RP12-L024
\item\label{refentry:REF202}B.Y. Levit, On the efficiency of a class of non-parametric estimates. Theory Probab. Appl. \textbf{20}(4), 723--740 (1975)\label{refend:REF202}
% REF203; 原书分章页 12; RP12-L025,RP12-L026,RP12-L027
\item\label{refentry:REF203}L. Li, E. Tchetgen Tchetgen, A.W. van der Vaart, J.M. Robins, Higher order inference on a treatment effect under low regularity conditions. Stat. Probab. Lett. \textbf{81}(7), 821--828 (2011)\label{refend:REF203}
% REF204; 原书分章页 12; RP12-L028,RP12-L029
\item\label{refentry:REF204}A. Liaw, M. Wiener, Classification and regression by randomforest. R News \textbf{2}(3), 18--22 (2002)\label{refend:REF204}
% REF205; 原书分章页 12; RP12-L030,RP12-L031,RP12-L032
\item\label{refentry:REF205}L. Liu, M.G. Hudgens, Large sample randomization inference of causal effects in the presence of interference. J. Am. Stat. Assoc. \textbf{109}(505), 288--301 (2014). ISSN 0162-1459\label{refend:REF205}
% REF206; 原书分章页 12; RP12-L033,RP12-L034
\item\label{refentry:REF206}Z. Liu, T. Stengos, Nonlinearities in cross country growth regressions: a semiparametric approach. J. Appl. Econom. \textbf{14}, 527--538 (1999)\label{refend:REF206}
% REF207; 原书分章页 12; RP12-L035,RP12-L036
\item\label{refentry:REF207}V. Loonis, X. Mary, Determinantal sampling designs. ArXiv e-prints, October (2015)\label{refend:REF207}
% REF208; 原书分章页 12; RP12-L037,RP12-L038
\item\label{refentry:REF208}B. Lu, R. Greevy, X. Xu, C. Beck, Optimal nonbipartite matching and its statistical applications. Am. Stat. \textbf{65}(1), 21--30 (2011)\label{refend:REF208}
% REF209; 原书分章页 12; RP12-L039,RP12-L040,RP12-L041
\item\label{refentry:REF209}A.R. Luedtke, M.J. van der Laan, Statistical inference for the mean outcome under a possibly non-unique optimal treatment strategy. Ann. Stat. \textbf{44}(2), 713--742 (2016a)\label{refend:REF209}
% REF210; 原书分章页 12; RP12-L042,RP12-L043
\item\label{refentry:REF210}A.R. Luedtke, M.J. van der Laan, Super-learning of an optimal dynamic treatment rule. Int. J. Biostat. \textbf{12}(1), 305--332 (2016b)\label{refend:REF210}
% REF211; 原书分章页 12; RP12-L044,RP12-L045
\item\label{refentry:REF211}A.R. Luedtke, M.J. van der Laan, Optimal individualized treatments in resource-limited settings. Int. J. Biostat. \textbf{12}(1), 283--303 (2016c)\label{refend:REF211}
% REF212; 原书分章页 13; RP13-L003,RP13-L004,RP13-L005
\item\label{refentry:REF212}A..R Luedtke, M. Carone, M.J. van der Laan, Discussion of deductive derivation and turing-computerization of semiparametric efficient estimation by Frangakis et al. Biometrics \textbf{71}(4), 875--879 (2015a)\label{refend:REF212}
% REF213; 原书分章页 13; RP13-L006,RP13-L007
\item\label{refentry:REF213}A.R. Luedtke, I. Díaz, M.J. van der Laan, The statistics of sensitivity analyses. Technical Report, Division of Biostatistics, University of California, Berkeley (2015b)\label{refend:REF213}
% REF214; 原书分章页 13; RP13-L008,RP13-L009
\item\label{refentry:REF214}K. Lum, Limitations of mitigating judicial bias with machine learning. Nat Hum. Behav. \textbf{1}, 0141 (2017)\label{refend:REF214}
% REF215; 原书分章页 13; RP13-L010,RP13-L011
\item\label{refentry:REF215}M. Lunn, D. McNeil, Applying Cox regression to competing risks. Biometrics \textbf{51}, 524--532 (1995). ISSN 0006-341X\label{refend:REF215}
% REF216; 原书分章页 13; RP13-L012,RP13-L013
\item\label{refentry:REF216}R. Lyons, Determinantal probability measures. Publications Mathématiques de l'Institut des Hautes Études Scientifiques \textbf{98}, 167--212 (2003)\label{refend:REF216}
% REF217; 原书分章页 13; RP13-L014,RP13-L015
\item\label{refentry:REF217}R. Lyons, The spread of evidence-poor medicine via flawed social-network analysis. Stat. Politics Policy \textbf{2}(1) 1--26 (2010)\label{refend:REF217}
% REF218; 原书分章页 13; RP13-L016,RP13-L017
\item\label{refentry:REF218}O. Macchi, The coincidence approach to stochastic point processes. Adv. Appl. Probab. \textbf{7}, 83--122 (1975)\label{refend:REF218}
% REF219; 原书分章页 13; RP13-L018,RP13-L019,RP13-L020
\item\label{refentry:REF219}R. Macklin, E. Cowan, Given financial constraints, it would be unethical to divert antiretroviral drugs from treatment to prevention. Health Aff. \textbf{31}(7), 1537--1544 (2012)\label{refend:REF219}
% REF220; 原书分章页 13; RP13-L021,RP13-L022,RP13-L023
\item\label{refentry:REF220}R.F. MacLehose, S. Kaufman, J.S. Kaufman, C. Poole, Bounding causal effects under uncontrolled confounding using counterfactuals. Epidemiology \textbf{16}(4), 548--555 (2005)\label{refend:REF220}
% REF221; 原书分章页 13; RP13-L024,RP13-L025
\item\label{refentry:REF221}E. Mammen, A.B. Tsybakov, Smooth discrimination analysis. Ann. Stat. \textbf{27}(6), 1808--1829 (1999)\label{refend:REF221}
% REF222; 原书分章页 13; RP13-L026,RP13-L027,RP13-L028,RP13-L029
\item\label{refentry:REF222}J.K. Mann, J.R. Balmes, T.A. Bruckner, K.M. Mortimer, H.G. Margolis, B. Pratt, S.K. Hammond, F.W. Lurmann, I.B. Tager, Short-term effects of air pollution on wheeze in asthmatic children in Fresno, California. Environ Health Perspect. \textbf{118}(10), 06 (2010)\label{refend:REF222}
% REF223; 原书分章页 13; RP13-L030,RP13-L031
\item\label{refentry:REF223}C.F. Manski, \textit{Partial} \textit{Identification} \textit{of} \textit{Probability} \textit{Distributions} (Springer, Berlin, Heidelberg, New York, 2003)\label{refend:REF223}
% REF224; 原书分章页 13; RP13-L032,RP13-L033
\item\label{refentry:REF224}C.F. Manski, Nonparametric bounds on treatment effects. Am. Econ. Rev. \textbf{80}, 319--323 (1990)\label{refend:REF224}
% REF225; 原书分章页 13; RP13-L034,RP13-L035
\item\label{refentry:REF225}D. Mayo, \textit{Error} \textit{and} \textit{the} \textit{Growth} \textit{of} \textit{Experimental} \textit{Knowledge} (University of Chicago Press, Chicago, 1996)\label{refend:REF225}
% REF226; 原书分章页 13; RP13-L036,RP13-L037,RP13-L038
\item\label{refentry:REF226}D. Mayo, \textit{Error} \textit{and} \textit{Inference:} \textit{Recent} \textit{Exchanges} \textit{on} \textit{Experimental} \textit{Reasoning,} \textit{Reliability,} \textit{and} \textit{the} \textit{Objectivity} \textit{and} \textit{Rationality} \textit{of} \textit{Science} (Cambridge, Chicago, 2010)\label{refend:REF226}
% REF227; 原书分章页 13; RP13-L039,RP13-L040
\item\label{refentry:REF227}S. Milborrow, T Hastie, R Tibshirani, Earth: multivariate adaptive regression spline models. R package version 3.2-7 (2014)\label{refend:REF227}
% REF228; 原书分章页 13; RP13-L041,RP13-L042,RP13-L043
\item\label{refentry:REF228}T. Mildenberger, Y. Rozenholc, D. Zasada, histogram: Construction of regular and irregular histograms with different options for automatic choice of bins (2009). \url{http://CRAN.R-project.org/package=histogram}\label{refend:REF228}
% REF229; 原书分章页 13; RP13-L044,RP13-L045
\item\label{refentry:REF229}E.E.M. Moodie, T.S. Richardson, D.A. Stephens, Demystifying optimal dynamic treatment regimes. Biometrics \textbf{63}(2), 447--455 (2007)\label{refend:REF229}
% REF230; 原书分章页 13,14; RP13-L046,RP13-L047,RP14-L003,RP14-L004
\item\label{refentry:REF230}K.L. Moore, M.J. van der Laan, Application of time-to-event methods in the assessment of safety in clinical trials, in \textit{Design,} \textit{Summarization,} \textit{Analysis} \textit{\&} \textit{Interpretation} \textit{of} \textit{Clinical} \textit{Trials} \textit{with} \textit{Time-to-Event} \textit{Endpoints}, ed. by K.E. Peace (Chapman \& Hall, Boca Raton, 2009a)\label{refend:REF230}
% REF231; 原书分章页 14; RP14-L005,RP14-L006,RP14-L007
\item\label{refentry:REF231}K.L. Moore, M.J. van der Laan, Covariate adjustment in randomized trials with binary outcomes: targeted maximum likelihood estimation. Stat. Med. \textbf{28}(1), 39--64 (2009b)\label{refend:REF231}
% REF232; 原书分章页 14; RP14-L008,RP14-L009,RP14-L010
\item\label{refentry:REF232}K.L. Moore, M.J. van der Laan, Increasing power in randomized trials with right censored outcomes through covariate adjustment. J. Biopharm. Stat. \textbf{19}(6), 1099--1131 (2009c)\label{refend:REF232}
% REF233; 原书分章页 14; RP14-L011,RP14-L012,RP14-L013
\item\label{refentry:REF233}K.L. Moore, R. Neugebauer, T. Valappil, M.J. van der Laan, Robust extraction of covariate information to improve estimation efficiency in randomized trials. Stat. Med. \textbf{30}(19), 2389--2408 (2011)\label{refend:REF233}
% REF234; 原书分章页 14; RP14-L014,RP14-L015
\item\label{refentry:REF234}N. Murata, A statistical study of on-line learning, in \textit{Online} \textit{Learning} \textit{and} \textit{Neural} \textit{Networks} (Cambridge University Press, Cambridge, 1998)\label{refend:REF234}
% REF235; 原书分章页 14; RP14-L016,RP14-L017
\item\label{refentry:REF235}S.A. Murphy, Optimal dynamic treatment regimes. J. R. Stat. Soc. Ser. B \textbf{65}(2), 331--66 (2003)\label{refend:REF235}
% REF236; 原书分章页 14; RP14-L018,RP14-L019
\item\label{refentry:REF236}S.A. Murphy, An experimental design for the development of adaptive treatment strategies. Stat. Med. \textbf{24}, 1455--1481 (2005)\label{refend:REF236}
% REF237; 原书分章页 14; RP14-L020,RP14-L021
\item\label{refentry:REF237}S.A. Murphy, M.J. van der Laan, J.M. Robins, Marginal mean models for dynamic treatment regimens. J. Am. Stat. Assoc. \textbf{96}, 1410--1424 (2001)\label{refend:REF237}
% REF238; 原书分章页 14; RP14-L022,RP14-L023
\item\label{refentry:REF238}E.A. Nadaraya, On estimating regression. Theory Probab. Appl. \textbf{9}(1), 141--142 (1964)\label{refend:REF238}
% REF239; 原书分章页 14; RP14-L024,RP14-L025,RP14-L026
\item\label{refentry:REF239}A.I Naimi, E.E.M. Moodie, N. Auger, J.S. Kaufman, Stochastic mediation contrasts in epidemiologic research: interpregnancy interval and the educational disparity in preterm delivery. Am. J. Epidemiol. \textbf{180}(4), 436--445 (2014)\label{refend:REF239}
% REF240; 原书分章页 14; RP14-L027,RP14-L028,RP14-L029,RP14-L030,RP14-L031
\item\label{refentry:REF240}D.M. Nathan, J.B. Buse, M.B. Davidson, E. Ferrannini, R.R. Holman, R. Sherwin, B. Zinman, Medical management of hyperglycemia in type 2 diabetes: a consensus algorithm for the initiation and adjustment of therapy: a consensus statement of the American Diabetes Association and the European Association for the Study of Diabetes. Diab. Care \textbf{32}(1), 193--203 (2009)\label{refend:REF240}
% REF241; 原书分章页 14; RP14-L032,RP14-L033,RP14-L034,RP14-L035,RP14-L036
\item\label{refentry:REF241}D.M. Nathan, P. A. Cleary, J.Y. Backlund, S.M. Genuth, J.M. Lachin, T.J. Orchard, P. Raskin, B. Zinman, Diabetes control and complications trial/epidemiology of diabetes interventions and complications (DCCT/EDIC) study research group. Intensive diabetes treatment and cardiovascular disease in patients with type 1 diabetes. N. Engl. J. Med. \textbf{22}(353), 2643--2653 (2005)\label{refend:REF241}
% REF242; 原书分章页 14; RP14-L037,RP14-L038,RP14-L039,RP14-L040,RP14-L041
\item\label{refentry:REF242}D.M. Nathan, J.B. Buse, M.B. Davidson, R.J. Heine, R.R. Holman, R. Sherwin, B. Zinman, Management of hyperglycemia in type 2 diabetes: a consensus algorithm for the initiation and adjustment of therapy: a consensus statement from the American Diabetes Association and the European Association for the Study of Diabetes. Diab. Care \textbf{29}, 1963--1972 (2006)\label{refend:REF242}
% REF243; 原书分章页 14; RP14-L042,RP14-L043
\item\label{refentry:REF243}NCEP (2002), Third Report of the National Cholesterol Education Program (NCEP) Expert Panel on Detection (2002)\label{refend:REF243}
% REF244; 原书分章页 14; RP14-L044,RP14-L045
\item\label{refentry:REF244}R. Neugebauer, J. Bullard, DSA: data-adaptive estimation with cross-validation and the D/S/A algorithm (2010). \url{http://www.stat.berkeley.edu/~laan/Software/}\label{refend:REF244}
% REF245; 原书分章页 14; RP14-L046,RP14-L047
\item\label{refentry:REF245}R. Neugebauer, M.J. van der Laan, Nonparametric causal effects based on marginal structural models. J. Stat. Plann. Infererence \textbf{137}(2), 419--434 (2007)\label{refend:REF245}
% REF246; 原书分章页 15; RP15-L003,RP15-L004,RP15-L005
\item\label{refentry:REF246}R. Neugebauer, M.J. Silverberg, M.J. van der Laan, Observational study and individualized antiretroviral therapy initiation rules for reducing cancer incidence in HIV-infected patients, chap. 26 (Springer, New York, 2011), pp. 436--456\label{refend:REF246}
% REF247; 原书分章页 15; RP15-L006,RP15-L007,RP15-L008,RP15-L009
\item\label{refentry:REF247}R. Neugebauer, B. Fireman, J.A. Roy, P.J. O'Connor, J.V. Selby, Dynamic marginal structural modeling to evaluate the comparative effectiveness of more or less aggressive treatment intensification strategies in adults with type 2 diabetes. Pharmacoepidemiol. Drug Saf. \textbf{21}(Suppl. 2), 99--113 (2012)\label{refend:REF247}
% REF248; 原书分章页 15; RP15-L010,RP15-L011,RP15-L012
\item\label{refentry:REF248}R. Neugebauer, B. Fireman, J.A. Roy, P.J. O'Connor, Impact of specific glucose-control strategies on microvascular and macrovascular outcomes in 58,000 adults with type 2 diabetes. Diab. Care \textbf{36}(11), 3510--3516 (2013)\label{refend:REF248}
% REF249; 原书分章页 15; RP15-L013,RP15-L014,RP15-L015
\item\label{refentry:REF249}R. Neugebauer, J. Schmittdiel, M.J. Laan, Targeted learning in real-world comparative effectiveness research with time-varying interventions. Stat. Med. \textbf{33}(14), 2480--2520 (2014a)\label{refend:REF249}
% REF250; 原书分章页 15; RP15-L016,RP15-L017,RP15-L018
\item\label{refentry:REF250}R. Neugebauer, J.A. Schmittdiel, Z. Zhu, J.A. Rassen, J.D. Seeger, S. Schneeweiss, High-dimensional propensity score algorithm in comparative effectiveness research with time-varying interventions. Stat. Med. \textbf{34}(5), 753--781 (2014b)\label{refend:REF250}
% REF251; 原书分章页 15; RP15-L019,RP15-L020,RP15-L021,RP15-L022
\item\label{refentry:REF251}R. Neugebauer, J.A. Schmittdiel, M.J. van der Laan, A case study of the impact of data-adaptive versus model-based estimation of the propensity scores on causal inferences from three inverse probability weighting estimators. Int. J. Biostat. \textbf{12}(1), 131--155 (2016)\label{refend:REF251}
% REF252; 原书分章页 15; RP15-L023,RP15-L024,RP15-L025
\item\label{refentry:REF252}J. Neyman, Sur les applications de la theorie des probabilites aux experiences agricoles: Essai des principes (In Polish). English translation by D.M. Dabrowska and T.P. Speed (1990). Stat. Sci. \textbf{5}, 465--480 (1923)\label{refend:REF252}
% REF253; 原书分章页 15; RP15-L026,RP15-L027,RP15-L028
\item\label{refentry:REF253}P.J. O'Connor, F. Ismail-Beigi, Near-normalization of glucose and microvascular diabetes complications: data from ACCORD and ADVANCE. Ther. Adv. Endocrinol. Metab. \textbf{2}(1), 17--26 (2011)\label{refend:REF253}
% REF254; 原书分章页 15; RP15-L029,RP15-L030
\item\label{refentry:REF254}E.L. Ogburn, T.J. VanderWeele, Vaccines, contagion, and social networks. ArXiv e-prints, March (2014)\label{refend:REF254}
% REF255; 原书分章页 15; RP15-L031,RP15-L032
\item\label{refentry:REF255}E.L. Ogburn, O. Sofrygin, M.J. van der Laan, I. Diaz, Causal inference for social network data with contagion. ArXiv e-prints, October (2017)\label{refend:REF255}
% REF256; 原书分章页 15; RP15-L033,RP15-L034
\item\label{refentry:REF256}B.A. Olken, Pre-analysis plans in economics. Technical report, Massachusetts Institute of Technology Department of Economics (2015)\label{refend:REF256}
% REF257; 原书分章页 15; RP15-L035,RP15-L036
\item\label{refentry:REF257}C. O'Neil, \textit{Weapons} \textit{of} \textit{Math} \textit{Destruction:} \textit{How} \textit{Big} \textit{Data} \textit{Increases} \textit{Inequality} \textit{and} \textit{Threatens} \textit{Democracy} (Crown Publishing Group, New York, 2016)\label{refend:REF257}
% REF258; 原书分章页 15; RP15-L037,RP15-L038,RP15-L039
\item\label{refentry:REF258}L. Orellana, A. Rotnitzky, J.M. Robins, Dynamic regime marginal structural mean models for estimation of optimal treatment regimes, part I: main content. Int. J. Biostat. \textbf{6}(2), Article 8 (2010)\label{refend:REF258}
% REF259; 原书分章页 15; RP15-L040,RP15-L041
\item\label{refentry:REF259}E. Parzen, On estimation of a probability density function and mode. Ann. Math. Stat. \textbf{33}(3), 1065--1076 (1962)\label{refend:REF259}
% REF260; 原书分章页 15; RP15-L042,RP15-L043,RP15-L044,RP15-L045
\item\label{refentry:REF260}A. Patel, S. MacMahon, J. Chalmers, B. Neal, L. Billot, M. Woodward, M. Marre, M. Cooper, P. Glasziou, D. Grobbee, P. Hamet, S. Harrap, S. Heller, Intensive blood glucose control and vascular outcomes in patients with type 2 diabetes. N. Engl. J. Med. \textbf{358}(24), 2560--2572 (2008)\label{refend:REF260}
% REF261; 原书分章页 15; RP15-L046
\item\label{refentry:REF261}J. Pearl, Causal diagrams for empirical research. Biometrika \textbf{82}, 669--710 (1995)\label{refend:REF261}
% REF262; 原书分章页 16; RP16-L003,RP16-L004
\item\label{refentry:REF262}J. Pearl, Direct and indirect effects, in \textit{Proceedings} \textit{of} \textit{the} \textit{17th} \textit{Conference} \textit{Uncertainty} \textit{in} \textit{Artificial} \textit{Intelligence} (Morgan Kaufmann, San Francisco, 2001)\label{refend:REF262}
% REF263; 原书分章页 16; RP16-L005,RP16-L006
\item\label{refentry:REF263}J. Pearl, \textit{Causality:} \textit{Models,} \textit{Reasoning,} \textit{and} \textit{Inference}, 2nd edn. (Cambridge, New York, 2009a)\label{refend:REF263}
% REF264; 原书分章页 16; RP16-L007,RP16-L008,RP16-L009
\item\label{refentry:REF264}J. Pearl, Myth, confusion, and science in causal analysis. Technical Report R-348, Cognitive Systems Laboratory, Computer Science Department University of California, Los Angeles, Los Angeles, CA, May 2009b\label{refend:REF264}
% REF265; 原书分章页 16; RP16-L010,RP16-L011
\item\label{refentry:REF265}J. Pearl, On the consistency rule in causal inference: axiom, definition, assumption, or theorem? Epidemiology \textbf{21}(6), 872--875 (2010)\label{refend:REF265}
% REF266; 原书分章页 16; RP16-L012,RP16-L013,RP16-L014
\item\label{refentry:REF266}J. Pearl, The mediation formula: a guide to the assessment of causal pathways in nonlinear models, in \textit{Causality:} \textit{Statistical} \textit{Perspectives} \textit{and} \textit{Applications}, ed. by C. Berzuini, P. Dawid, L. Bernardinelli (Springer, Berlin, 2011)\label{refend:REF266}
% REF267; 原书分章页 16; RP16-L015,RP16-L016
\item\label{refentry:REF267}R. Pemantle, Y. Peres, Concentration of Lipschitz functionals of determinantal and other strong Rayleigh measures. Comb. Probab. Comput. \textbf{23}(1), 140--160 (2014)\label{refend:REF267}
% REF268; 原书分章页 16; RP16-L017,RP16-L018,RP16-L019
\item\label{refentry:REF268}W.D. Penny, K.E. Stephan, A. Mechelli, K.J. Friston, Modelling functional integration: a comparison of structural equation and dynamic causal models. Neuroimage \textbf{23}, S264--S274 (2004)\label{refend:REF268}
% REF269; 原书分章页 16; RP16-L020,RP16-L021
\item\label{refentry:REF269}G. Peoples, New study from Pandora touts the Pandora effect on music sales. Billboard (2014)\label{refend:REF269}
% REF270; 原书分章页 16; RP16-L022,RP16-L023
\item\label{refentry:REF270}A. Peters, T. Hothorn, ipred: improved predictors (2009) \url{http://CRAN.R-project.org/package=ipred}\label{refend:REF270}
% REF271; 原书分章页 16; RP16-L024,RP16-L025,RP16-L026
\item\label{refentry:REF271}M. Petersen, J. Schwab, S. Gruber, N. Blaser, M. Schomaker, M.J. van der Laan, Targeted maximum likelihood estimation for dynamic and static longitudinal marginal structural working models. J. Causal Inference \textbf{2}(2), 147--185 (2014)\label{refend:REF271}
% REF272; 原书分章页 16; RP16-L027,RP16-L028,RP16-L029,RP16-L030
\item\label{refentry:REF272}M.L. Petersen, E. LeDell, J. Schwab, V. Sarovar, R. Gross, N. Reynolds, J.E. Haberer, K. Goggin, C. Golin, J. Arnsten et al., Super learner analysis of electronic adherence data improves viral prediction and may provide strategies for selective HIV RNA monitoring. J. Acquir. Immune Defic. Syndr. \textbf{69}(1), 109 (2015)\label{refend:REF272}
% REF273; 原书分章页 16; RP16-L031,RP16-L032
\item\label{refentry:REF273}J. Pfanzagl, \textit{Contributions} \textit{to} \textit{a} \textit{General} \textit{Asymptotic} \textit{Statistical} \textit{Theory} (Springer, Berlin, 1982)\label{refend:REF273}
% REF274; 原书分章页 16; RP16-L033,RP16-L034
\item\label{refentry:REF274}J. Pfanzagl, \textit{Asymptotic} \textit{Expansions} \textit{for} \textit{General} \textit{Statistical} \textit{Models}, vol. 31 (Springer, Berlin, 1985)\label{refend:REF274}
% REF275; 原书分章页 16; RP16-L035,RP16-L036
\item\label{refentry:REF275}J. Pfanzagl, \textit{Estimation} \textit{in} \textit{Semiparametric} \textit{Models} (Springer, Berlin, Heidelberg, New York, 1990)\label{refend:REF275}
% REF276; 原书分章页 16; RP16-L037,RP16-L038
\item\label{refentry:REF276}I. Phyllis, F. Russo. \textit{Causality;} \textit{Philosophical} \textit{Theory} \textit{meets} \textit{Scientific} \textit{Practice} (Oxford University Press, Oxford, 2016)\label{refend:REF276}
% REF277; 原书分章页 16; RP16-L039,RP16-L040
\item\label{refentry:REF277}M. Pintilie, Analysing and interpreting competing risk data. Stat. Med. \textbf{26}(6), 1360--1367 (2007)\label{refend:REF277}
% REF278; 原书分章页 16; RP16-L041,RP16-L042,RP16-L043
\item\label{refentry:REF278}R. Pirracchio, M.L. Petersen, M.J. van der Laan, Improving propensity score estimators' robustness to model misspecification using super learner. Am. J. Epidemiol. \textbf{181}(2), 108--119 (2014)\label{refend:REF278}
% REF279; 原书分章页 16; RP16-L044,RP16-L045,RP16-L046,RP16-L047
\item\label{refentry:REF279}R. Pirracchio, M.L. Petersen, M. Carone, M.R. Rigon, S. Chevret, M.J. van der Laan, Mortality prediction in intensive care units with the super ICU learner algorithm (SICULA): a population-based study. Lancet Respir. Med. \textbf{3}(1), 42--52 (2015)\label{refend:REF279}
% REF280; 原书分章页 17; RP17-L003,RP17-L004
\item\label{refentry:REF280}R.W. Platt, E.F. Schisterman, S.R. Cole, Time-modified confounding. Am. J. Epidemiol. \textbf{170}(6), 687--694 (2009)\label{refend:REF280}
% REF281; 原书分章页 17; RP17-L005,RP17-L006,RP17-L007
\item\label{refentry:REF281}S.J. Pocock, S.E. Assmann, L.E. Enos, L.E. Kasten, Subgroup analysis, covariate adjustment, and baseline comparisons in clinical trial reporting: current practice and problems. Stat. Med. \textbf{21}, 2917--2930 (2002)\label{refend:REF281}
% REF282; 原书分章页 17; RP17-L008,RP17-L009
\item\label{refentry:REF282}E.C. Polley, M.J. van der Laan, SuperLearner: super learner prediction (2013). \url{http://CRAN.R-project.org/package=SuperLearner}\label{refend:REF282}
% REF283; 原书分章页 17; RP17-L010,RP17-L011,RP17-L012,RP17-L013
\item\label{refentry:REF283}E.C. Polley, M.J. van der Laan, Predicting optimal treatment assignment based on prognostic factors in cancer patients. in \textit{Design,} \textit{Summarization,} \textit{Analysis} \textit{\&} \textit{Interpretation} \textit{of} \textit{Clinical} \textit{Trials} \textit{with} \textit{Time-to-Event} \textit{Endpoints}, ed. by K.E. Peace (Boca Raton, Chapman \& Hall, 2009)\label{refend:REF283}
% REF284; 原书分章页 17; RP17-L014,RP17-L015
\item\label{refentry:REF284}E.C. Polley, M.J. van der Laan, Super learner in prediction. Technical Report 266, Division of Biostatistics, University of California, Berkeley (2010)\label{refend:REF284}
% REF285; 原书分章页 17; RP17-L016,RP17-L017,RP17-L018
\item\label{refentry:REF285}E.C Polley, S. Rose, M.J. van der Laan, Super-learning, in \textit{Targeted} \textit{Learning:} \textit{Causal} \textit{Inference} \textit{for} \textit{Observational} \textit{and} \textit{Experimental} \textit{Data}, ed. by M.J. van der Laan, S. Rose (Springer, Berlin, Heidelberg, New York, 2011)\label{refend:REF285}
% REF286; 原书分章页 17; RP17-L019,RP17-L020
\item\label{refentry:REF286}E.C. Polley, E. LeDell, C. Kennedy, M.J. van der Laan, SuperLearner: super learner prediction (2017). \url{https://github.com/ecpolley/SuperLearner}\label{refend:REF286}
% REF287; 原书分章页 17; RP17-L021,RP17-L022
\item\label{refentry:REF287}B.T. Polyak, A.B. Juditsky, Acceleration of stochastic approximation by averaging. SIAM J. Control. Optim. \textbf{30}(4), 838--855 (1992)\label{refend:REF287}
% REF288; 原书分章页 17; RP17-L023,RP17-L024
\item\label{refentry:REF288}T.M. Porter, \textit{The} \textit{Rise} \textit{of} \textit{Statistical} \textit{Thinking} (Princeton University Press, Princeton, 1986)\label{refend:REF288}
% REF289; 原书分章页 17; RP17-L025,RP17-L026
\item\label{refentry:REF289}T.M. Porter, \textit{Trust} \textit{in} \textit{Numbers:} \textit{The} \textit{Pursuit} \textit{of} \textit{Objectivity} \textit{in} \textit{Science} \textit{and} \textit{Public} \textit{Life} (Princeton University Press, Princeton, 1995)\label{refend:REF289}
% REF290; 原书分章页 17; RP17-L027,RP17-L028
\item\label{refentry:REF290}K.E Porter, S. Gruber, M.J. van der Laan, J.S. Sekhon, The relative performance of targeted maximum likelihood estimators. Int. J. Biostat. 7(1) (2011)\label{refend:REF290}
% REF291; 原书分章页 17; RP17-L029,RP17-L030,RP17-L031
\item\label{refentry:REF291}R.L. Prentice, J.D. Kalbfleisch, A.V. Peterson Jr, N. Flournoy, V.T. Farewell, N.E. Breslow, The analysis of failure times in the presence of competing risks. Biometrics \textbf{34}(4), 541--554 (1978)\label{refend:REF291}
% REF292; 原书分章页 17; RP17-L032,RP17-L033
\item\label{refentry:REF292}M. Qian, S.A. Murphy, Performance guarantees for individualized treatment rules. Ann. Stat. \textbf{39}(2), 1180--1210 (2011)\label{refend:REF292}
% REF293; 原书分章页 17; RP17-L034,RP17-L035,RP17-L036
\item\label{refentry:REF293}R Development Core Team. R: A language and environment for statistical computing. R Foundation for Statistical Computing, Vienna (2016). \url{http://www.R-project.org}.\label{refend:REF293}
% REF294; 原书分章页 17; RP17-L037,RP17-L038,RP17-L039,RP17-L040
\item\label{refentry:REF294}K.K. Ray, S.R. Seshasai, S. Wijesuriya, R. Sivakumaran, S. Nethercott, D. Preiss, S. Erqou, N. Sattar, Effect of intensive control of glucose on cardiovascular outcomes and death in patients with diabetes mellitus: a meta-analysis of randomised controlled trials. Lancet \textbf{373}, 1765--72 (2009)\label{refend:REF294}
% REF295; 原书分章页 17; RP17-L041,RP17-L042,RP17-L043
\item\label{refentry:REF295}J.M. Robins, A new approach to causal inference in mortality studies with sustained exposure periods--application to control of the healthy worker survivor effect. Math. Modell. \textbf{7}, 1393--1512 (1986)\label{refend:REF295}
% REF296; 原书分章页 17; RP17-L044,RP17-L045,RP17-L046
\item\label{refentry:REF296}J.M. Robins, Addendum to: ``A new approach to causal inference in mortality studies with a sustained exposure period---application to control of the healthy worker survivor effect''. Comput. Math. Appl. \textbf{14}(9--12), 923--945 (1987)\label{refend:REF296}
% REF297; 原书分章页 18; RP18-L003,RP18-L004
\item\label{refentry:REF297}J.M. Robins, Marginal structural models, in \textit{1997} \textit{Proceedings} \textit{of} \textit{the} \textit{American} \textit{Statistical} \textit{Association.} \textit{Section} \textit{on} \textit{Bayesian} \textit{Statistical} \textit{Science}, pp. 1--10 (1998)\label{refend:REF297}
% REF298; 原书分章页 18; RP18-L005,RP18-L006
\item\label{refentry:REF298}J.M. Robins, Association, causation and marginal structural models. Synthese \textbf{121}, 151--179 (1999)\label{refend:REF298}
% REF299; 原书分章页 18; RP18-L007,RP18-L008
\item\label{refentry:REF299}J.M. Robins, Robust estimation in sequentially ignorable missing data and causal inference models, in \textit{Proceedings} \textit{of} \textit{the} \textit{American} \textit{Statistical} \textit{Association} (2000)\label{refend:REF299}
% REF300; 原书分章页 18; RP18-L009,RP18-L010,RP18-L011
\item\label{refentry:REF300}J.M. Robins, Optimal structural nested models for optimal sequential decisions, in \textit{Proceedings} \textit{of} \textit{the} \textit{Second} \textit{Seattle} \textit{Symposium} \textit{in} \textit{Biostatistics:} \textit{Analysis} \textit{of} \textit{Correlated} \textit{Data} (2004)\label{refend:REF300}
% REF301; 原书分章页 18; RP18-L012,RP18-L013
\item\label{refentry:REF301}J.M. Robins, S. Greenland, Identifiability and exchangeability for direct and indirect effects. Epidemiol \textbf{3}, 143--155 (1992)\label{refend:REF301}
% REF302; 原书分章页 18; RP18-L014,RP18-L015
\item\label{refentry:REF302}J.M. Robins, Y. Ritov, Toward a curse of dimensionality appropriate (coda) asymptotic theory for semi-parametric models. Stat. Med. \textbf{16}, 285--319 (1997)\label{refend:REF302}
% REF303; 原书分章页 18; RP18-L016,RP18-L017,RP18-L018
\item\label{refentry:REF303}J.M. Robins, A. Rotnitzky, Recovery of information and adjustment for dependent censoring using surrogate markers, in \textit{AIDS} \textit{Epidemiology} (Birkhäuser, Basel, 1992)\label{refend:REF303}
% REF304; 原书分章页 18; RP18-L019,RP18-L020,RP18-L021
\item\label{refentry:REF304}J.M. Robins, A. Rotnitzky, L.P. Zhao, Estimation of regression coefficients when some regressors are not always observed. J. Am. Stat. Assoc. \textbf{89}(427), 846--866 (1994)\label{refend:REF304}
% REF305; 原书分章页 18; RP18-L022,RP18-L023,RP18-L024,RP18-L025
\item\label{refentry:REF305}J.M. Robins, A. Rotnitzky, D.O. Scharfstein, Sensitivity analysis for selection bias and unmeasured confounding in missing data and causal inference models, in \textit{Statistical} \textit{Models} \textit{in} \textit{Epidemiology,} \textit{the} \textit{Environment} \textit{and} \textit{Clinical} \textit{Trials}. IMA Volumes in Mathematics and Its Applications (Springer, Berlin, 1999)\label{refend:REF305}
% REF306; 原书分章页 18; RP18-L026,RP18-L027
\item\label{refentry:REF306}J.M. Robins, M.A. Hernan, B. Brumback, Marginal structural models and causal inference in epidemiology. Epidemiology \textbf{11}(5), 550--560 (2000)\label{refend:REF306}
% REF307; 原书分章页 18; RP18-L028,RP18-L029,RP18-L030,RP18-L031
\item\label{refentry:REF307}J.M. Robins, M.A. Hernán, U. Siebert, Effects of multiple interventions, in \textit{Comparative} \textit{Quantification} \textit{of} \textit{Health} \textit{Risks:} \textit{Global} \textit{and} \textit{Regional} \textit{Burden} \textit{of} \textit{Disease} \textit{Attributable} \textit{to} \textit{Selected} \textit{Major} \textit{Risk} \textit{Factors}, vol. 1 (World Health Organization, Geneva, 2004), pp. 2191--2230\label{refend:REF307}
% REF308; 原书分章页 18; RP18-L032,RP18-L033,RP18-L034,RP18-L035
\item\label{refentry:REF308}J.M. Robins, L. Li, E. Tchetgen Tchetgen, A.W. van der Vaart, Higher order influence functions and minimax estimation of nonlinear functionals, in \textit{Probability} \textit{and} \textit{Statistics:} \textit{Essays} \textit{in} \textit{Honor} \textit{of} \textit{David} \textit{A.} \textit{Freedman}, (Institute of Mathematical Statistics, 2008a), pp. 335--421\label{refend:REF308}
% REF309; 原书分章页 18; RP18-L036,RP18-L037
\item\label{refentry:REF309}J.M. Robins, L. Orellana, A. Rotnitzky, Estimation and extrapolation of optimal treatment and testing strategies. Stat. Med. \textbf{27}, 4678--4721 (2008b)\label{refend:REF309}
% REF310; 原书分章页 18; RP18-L038,RP18-L039
\item\label{refentry:REF310}J.M. Robins, L. Li, E. Tchetgen Tchetgen, A.W. van der Vaart, Quadratic Semiparametric Von Mises calculus. Metrika \textbf{69}(2--3), 227--247 (2009)\label{refend:REF310}
% REF311; 原书分章页 18; RP18-L040,RP18-L041,RP18-L042,RP18-L043
\item\label{refentry:REF311}M. Rolland, P.T. Edlefsen, B.B. Larsen, S. Tovanabutra, E. Sanders-Buell, T. Hertz, C. Carrico, S. Menis, C.A. Magaret, H. Ahmed, Increased HIV-1 vaccine efficacy against viruses with genetic signatures in Env V2. Nature \textbf{490}(7420), 417--420 (2012). ISSN 0028-0836\label{refend:REF311}
% REF312; 原书分章页 18; RP18-L044,RP18-L045
\item\label{refentry:REF312}S. Rose, Mortality risk score prediction in an elderly population using machine learning. Am. J. Epidemiol. \textbf{177}(5), 443--452 (2013)\label{refend:REF312}
% REF313; 原书分章页 18; RP18-L046,RP18-L047
\item\label{refentry:REF313}S. Rose, Targeted learning for pre-analysis plans in public health and health policy research. Observational Stud. \textbf{1}, 294--306 (2015)\label{refend:REF313}
% REF314; 原书分章页 19; RP19-L003,RP19-L004
\item\label{refentry:REF314}S. Rose, A machine learning framework for plan payment risk adjustment. Health Serv. Res. \textbf{51}(6), 2358--2374 (2016)\label{refend:REF314}
% REF315; 原书分章页 19; RP19-L005,RP19-L006
\item\label{refentry:REF315}S. Rose, Robust machine learning variable importance analyses of medical conditions for health care spending. Health Serv. Res. (2018, in press)\label{refend:REF315}
% REF316; 原书分章页 19; RP19-L007,RP19-L008
\item\label{refentry:REF316}S. Rose, S. Bergquist, T. Layton, Computational health economics for identification of unprofitable health care enrollees. Biostatistics \textbf{18}(4), 682--694 (2017)\label{refend:REF316}
% REF317; 原书分章页 19; RP19-L009,RP19-L010
\item\label{refentry:REF317}S. Rose, M.J. van der Laan, Simple optimal weighting of cases and controls in case-control studies. Int. J. Biostat. \textbf{4}(1), Article 19 (2008)\label{refend:REF317}
% REF318; 原书分章页 19; RP19-L011,RP19-L012
\item\label{refentry:REF318}S. Rose, M.J. van der Laan, Why match? Investigating matched case-control study designs with causal effect estimation. Int. J. Biostat. \textbf{5}(1), Article 1 (2009)\label{refend:REF318}
% REF319; 原书分章页 19; RP19-L013,RP19-L014
\item\label{refentry:REF319}S. Rose, M.J. van der Laan, A targeted maximum likelihood estimator for two-stage designs. Int. J. Biostat. \textbf{7}(1), Article 17 (2011)\label{refend:REF319}
% REF320; 原书分章页 19; RP19-L015,RP19-L016
\item\label{refentry:REF320}S. Rose, M.J. van der Laan, A double robust approach to causal effects in case-control studies. Am. J. Epidemiol. \textbf{179}(6), 663--669 (2014a)\label{refend:REF320}
% REF321; 原书分章页 19; RP19-L017,RP19-L018
\item\label{refentry:REF321}S. Rose, M.J. van der Laan, Rose and van der Laan respond to ``Some advantages of RERI''. Am. J. Epidemiol. \textbf{179}(6), 672--673 (2014b)\label{refend:REF321}
% REF322; 原书分章页 19; RP19-L019,RP19-L020,RP19-L021
\item\label{refentry:REF322}P.R. Rosenbaum, D.B. Rubin, Assessing sensitivity to an unobserved binary covariate in an observational study with binary outcome. J. R. Stat. Soc. Ser. B \textbf{45}, 212--218 (1983a)\label{refend:REF322}
% REF323; 原书分章页 19; RP19-L022,RP19-L023
\item\label{refentry:REF323}P.R. Rosenbaum, Interference Between Units in Randomized Experiments. J. Am. Stat. Assoc. \textbf{102}(477), 191--200 (2007)\label{refend:REF323}
% REF324; 原书分章页 19; RP19-L024,RP19-L025
\item\label{refentry:REF324}P.R. Rosenbaum, D.B. Rubin, The central role of the propensity score in observational studies for causal effects. Biometrika \textbf{70}, 41--55 (1983b)\label{refend:REF324}
% REF325; 原书分章页 19; RP19-L026,RP19-L027
\item\label{refentry:REF325}M. Rosenblatt, Remarks on some nonparametric estimates of a density function. Ann. Math. Stat. \textbf{27}(3), 832--837 (1956)\label{refend:REF325}
% REF326; 原书分章页 19; RP19-L028,RP19-L029,RP19-L030
\item\label{refentry:REF326}M. Rosenblum, M.J. van der Laan, Using regression models to analyze randomized trials: asymptotically valid hypothesis tests despite incorrectly specified models. Biometrics \textbf{65}(3), 937--945 (2009)\label{refend:REF326}
% REF327; 原书分章页 19; RP19-L031,RP19-L032
\item\label{refentry:REF327}M. Rosenblum, M.J. van der Laan, Targeted maximum likelihood estimation of the parameter of a marginal structural model. Int. J. Biostat. \textbf{6}(2), 19 (2010a)\label{refend:REF327}
% REF328; 原书分章页 19; RP19-L033,RP19-L034,RP19-L035
\item\label{refentry:REF328}M. Rosenblum, M.J. van der Laan, Simple, efficient estimators of treatment effects in randomized trials using generalized linear models to leverage baseline variables. Int. J. Biostat. \textbf{6}(1), Article 13 (2010b)\label{refend:REF328}
% REF329; 原书分章页 19; RP19-L036,RP19-L037,RP19-L038
\item\label{refentry:REF329}M. Rosenblum, S.G. Deeks, M.J. van der Laan, D.R. Bangsberg, The risk of virologic failure decreases with duration of HIV suppression, at greater than 50\% adherence to antiretroviral therapy. PLoS ONE \textbf{4}(9), e7196 (2009)\label{refend:REF329}
% REF330; 原书分章页 19; RP19-L039,RP19-L040,RP19-L041
\item\label{refentry:REF330}R.H. Rosenman, M. Friedman, R. Straus, M. Wurm, R. Kositchek, W. Hahn, N.T. Werthessen, A predictive study of coronary heart disease: the western collaborative group study. J. Am. Med. Assoc. \textbf{189}(1), 15--22 (1964)\label{refend:REF330}
% REF331; 原书分章页 19; RP19-L042,RP19-L043,RP19-L044
\item\label{refentry:REF331}R.H. Rosenman, R.J. Brand, C.D. Jenkins, M. Friedman, R. Straus, M. Wurm, Coronary heart disease in the western collaborative group study: final follow-up experience of 8 1/2 years. J. Am. Med. Assoc. \textbf{233}(8), 872--877 (1975)\label{refend:REF331}
% REF332; 原书分章页 19; RP19-L045
\item\label{refentry:REF332}B. Rosner, \textit{Fundamentals} \textit{of} \textit{Biostatistics}, 5th edn. (Duxbury, Pacific Grove, 1999)\label{refend:REF332}
% REF333; 原书分章页 20; RP20-L003,RP20-L004,RP20-L005
\item\label{refentry:REF333}S. Rosthø j, C. Fullwood, R. Henderson, S. Stewart, Estimation of optimal dynamic anticoagulation regimes from observational data: a regret-based approach. Stat. Med. \textbf{88}, 4197--4215 (2006)\label{refend:REF333}
% REF334; 原书分章页 20; RP20-L006,RP20-L007,RP20-L008
\item\label{refentry:REF334}A. Rotnitzky, D. Scharfstein, S. Ting-Li Su, J. Robins, Methods for conducting sensitivity analysis of trials with potentially nonignorable competing causes of censoring. Biometrics \textbf{57}(1), 103--113 (2001)\label{refend:REF334}
% REF335; 原书分章页 20; RP20-L009,RP20-L010,RP20-L011
\item\label{refentry:REF335}A. Rotnitzky, J.M. Robins, D.O. Scharfstein, Semiparametric regression for repeated outcomes with nonignorable nonresponse. J. Am. Med. Assoc. \textbf{93}(444), 1321--1339 (1998)\label{refend:REF335}
% REF336; 原书分章页 20; RP20-L012,RP20-L013,RP20-L014
\item\label{refentry:REF336}Y. Rozenholc, T. Mildenberger, U. Gather, Combining regular and irregular histograms by penalized likelihood. Comput. Stat. Data Anal. \textbf{54}(12), 3313--3323 (2010)\label{refend:REF336}
% REF337; 原书分章页 20; RP20-L015,RP20-L016
\item\label{refentry:REF337}D.B. Rubin, Randomization analysis of experimental data: The fisher randomization test comment. J. Am. Stat. Assoc. \textbf{75}(371), 591--593 (1980)\label{refend:REF337}
% REF338; 原书分章页 20; RP20-L017,RP20-L018
\item\label{refentry:REF338}D.B. Rubin, Estimating causal effects of treatments in randomized and nonrandomized studies. J. Educ. Psychol. \textbf{66}, 688--701 (1974)\label{refend:REF338}
% REF339; 原书分章页 20; RP20-L019,RP20-L020,RP20-L021
\item\label{refentry:REF339}D.B. Rubin, Multivariate matching methods that are equal percent bias reducing, II: maximums on bias reduction for fixed sample sizes. Biometrics \textbf{32}(1), 121--132 (1976)\label{refend:REF339}
% REF340; 原书分章页 20; RP20-L022,RP20-L023
\item\label{refentry:REF340}D.B. Rubin, Comment: Neyman (1923) and causal inference in experiments and observational studies. Stat. Sci. \textbf{5}(4), 472--480 (1990)\label{refend:REF340}
% REF341; 原书分章页 20; RP20-L024,RP20-L025
\item\label{refentry:REF341}D.B. Rubin, \textit{Matched} \textit{Sampling} \textit{for} \textit{Causal} \textit{Effects} (Cambridge, Cambridge, MA, 2006)\label{refend:REF341}
% REF342; 原书分章页 20; RP20-L026,RP20-L027,RP20-L028
\item\label{refentry:REF342}D.B. Rubin, M.J. van der Laan, Empirical efficiency maximization: improved locally efficient covariate adjustment in randomized experiments and survival analysis. Int. J. Biostat. \textbf{4}(1), Article 5 (2008)\label{refend:REF342}
% REF343; 原书分章页 20; RP20-L029,RP20-L030
\item\label{refentry:REF343}D.B. Rubin, M.J. van der Laan, Targeted ANCOVA estimator in RCTs, in \textit{Targeted} \textit{Learning} (Springer, Berlin, 2011), pp. 201--215\label{refend:REF343}
% REF344; 原书分章页 20; RP20-L031,RP20-L032
\item\label{refentry:REF344}D.B. Rubin, M.J. van der Laan, Statistical issues and limitations in personalized medicine research with clinical trials. Int. J. Biostat. \textbf{8}(1), Article 1 (2012)\label{refend:REF344}
% REF345; 原书分章页 20; RP20-L033,RP20-L034
\item\label{refentry:REF345}M. Rudelson, R. Vershynin, Hanson-Wright inequality and subGaussian concentration. Electron. Commun. Probab. \textbf{18}(82), 1--9 (2013)\label{refend:REF345}
% REF346; 原书分章页 20; RP20-L035,RP20-L036
\item\label{refentry:REF346}M.R. Sampford, On sampling without replacement with unequal probabilities of selection. Biometrika \textbf{54}(3--4), 499--513 (1967)\label{refend:REF346}
% REF347; 原书分章页 20; RP20-L037,RP20-L038,RP20-L039
\item\label{refentry:REF347}S. Sapp, M.J. van der Laan, K. Page, Targeted estimation of binary variable importance measures with interval-censored outcomes. Int. J. Biostat. \textbf{10}(1), 77--97 (2014)\label{refend:REF347}
% REF348; 原书分章页 20; RP20-L040,RP20-L041
\item\label{refentry:REF348}D.O. Scharfstein, J.M. Robins, Estimation of the failure time distribution in the presence of informative censoring. Biometrika \textbf{89}(3), 617--634 (2002)\label{refend:REF348}
% REF349; 原书分章页 20; RP20-L042,RP20-L043,RP20-L044
\item\label{refentry:REF349}D.O. Scharfstein, A. Rotnitzky, J.M. Robins, Adjusting for nonignorable drop-out using semiparametric nonresponse models, (with discussion and rejoinder). J. Am. Stat. Assoc. \textbf{94}, 1096--1120, 1121--1146 (1999)\label{refend:REF349}
% REF350; 原书分章页 20; RP20-L045,RP20-L046,RP20-L047
\item\label{refentry:REF350}M.E. Schnitzer, J. Lok, S. Gruber, Variable selection for confounder control, flexible modeling and collaborative targeted minimum loss-based estimation in causal inference. Int. J. Biostat. \textbf{12}(1), 97--115 (2016)\label{refend:REF350}
% REF351; 原书分章页 21; RP21-L003,RP21-L004,RP21-L005
\item\label{refentry:REF351}M.E. Schnitzer, M.J. van der Laan, E.E.M. Moodie, R.W. Platt, Effect of breast-feeding on gastrointestinal infection in infants: a targeted maximum likelihood approach for clustered longitudinal data. Ann. Appl. Stat. \textbf{8}(2), 703--725 (2014)\label{refend:REF351}
% REF352; 原书分章页 21; RP21-L006,RP21-L007
\item\label{refentry:REF352}P. Schochet, Estimators for clustered education RCTs using the Neyman model for causal inference. J. Educ. Behav. Stat. \textbf{38}(3), 219--238 (2013)\label{refend:REF352}
% REF353; 原书分章页 21; RP21-L008,RP21-L009
\item\label{refentry:REF353}M.S. Schuler, S. Rose, Targeted maximum likelihood estimation for causal inference in observational studies. Am. J. Epidemiol. \textbf{185}(1), 65--73 (2017)\label{refend:REF353}
% REF354; 原书分章页 21; RP21-L010,RP21-L011
\item\label{refentry:REF354}S. Selvaraj, V. Prasad. Characteristics of cluster randomized trials: Are they living up to the randomized trial? JAMA Intern. Med. \textbf{173}(23), 313 (2013)\label{refend:REF354}
% REF355; 原书分章页 21; RP21-L012,RP21-L013
\item\label{refentry:REF355}S. Shalev-Shwartz, Online learning and online convex optimization. Found. Trends Mach. Learn. \textbf{4}(2), 107--194 (2011)\label{refend:REF355}
% REF356; 原书分章页 21; RP21-L014,RP21-L015
\item\label{refentry:REF356}S. Shalev-Shwartz, Y. Singer, N. Srebro, A. Cotter, Pegasos: primal estimated sub-gradient solver for SVM. Math. Programm. \textbf{127}(1), 3--30 (2011)\label{refend:REF356}
% REF357; 原书分章页 21; RP21-L016,RP21-L017,RP21-L018
\item\label{refentry:REF357}C.R. Shalizi, A.C. Thomas, Homophily and contagion are generically confounded in observational social network studies. Sociol. Methods Res. \textbf{40}(2), 211--239 (2011)\label{refend:REF357}
% REF358; 原书分章页 21; RP21-L019,RP21-L020
\item\label{refentry:REF358}C. Shen, X. Li, L. Li, Inverse probability weighting for covariate adjustment in randomized studies. Stat. Med. \textbf{33}, 555--568 (2014)\label{refend:REF358}
% REF359; 原书分章页 21; RP21-L021,RP21-L022
\item\label{refentry:REF359}A. Shrestha, S. Bergquist, E. Montz, S. Rose, Mental health risk adjustment with clinical categories and machine learning. Health Serv. Res. (2018, in press)\label{refend:REF359}
% REF360; 原书分章页 21; RP21-L023,RP21-L024
\item\label{refentry:REF360}J.A. Singh, Antiretroviral resource allocation for HIV prevention. AIDS \textbf{27}(6), 863--865 (2013)\label{refend:REF360}
% REF361; 原书分章页 21; RP21-L025,RP21-L026
\item\label{refentry:REF361}S.E. Sinisi, M.J. van der Laan, Deletion/Substitution/Addition algorithm in learning with applications in genomics. Stat. Appl. Genet. Mol. \textbf{3}(1), Article 18 (2004)\label{refend:REF361}
% REF362; 原书分章页 21; RP21-L027,RP21-L028,RP21-L029,RP21-L030,RP21-L031,RP21-L032,RP21-L033
\item\label{refentry:REF362}J.S. Skyler, R. Bergenstal, R.O. Bonow, J. Buse, P. Deedwania, E.A.M. Gale, B.V. Howard, M.S. Kirkman, M. Kosiborod, P. Reaven, R.S. Sherwin, Intensive Glycemic Control and the prevention of cardiovascular events: implications of the accord, advance, and VA diabetes trials: a position statement of the American Diabetes Association and a scientific statement of the American College of Cardiology Foundation and the American Heart Association. Diab. Care \textbf{32}, 187--92 (2009)\label{refend:REF362}
% REF363; 原书分章页 21; RP21-L034,RP21-L035,RP21-L036,RP21-L037
\item\label{refentry:REF363}J.W. Smith, J.E. Everhart, W.C. Dickson, W.C. Knowler, R.S. Johannes, Using the adap learning algorithm to forecast the onset of diabetes mellitus, in \textit{Proceedings} \textit{of} \textit{the} \textit{Annual} \textit{Symposium} \textit{on} \textit{Computer} \textit{Application} \textit{in} \textit{Medical} \textit{Care} (American Medical Informatics Association, Bethesda, 1988), p. 261\label{refend:REF363}
% REF364; 原书分章页 21; RP21-L038,RP21-L039,RP21-L040
\item\label{refentry:REF364}J.M. Snowden, S. Rose, K.M. Mortimer, Implementation of g-computation on a simulated data set: demonstration of a causal inference technique. Am. J. Epidemiol. \textbf{173}(7), 731--738 (2011)\label{refend:REF364}
% REF365; 原书分章页 21; RP21-L041,RP21-L042
\item\label{refentry:REF365}M. Sobel, What do randomized studies of housing mobility demonstrate? J. Am. Stat. Assoc. \textbf{101}(476), 1398--1407 (2006)\label{refend:REF365}
% REF366; 原书分章页 21; RP21-L043,RP21-L044,RP21-L045
\item\label{refentry:REF366}O. Sofrygin, M.J. van der Laan, R. Neugebauer, Simcausal R package: conducting transparent and reproducible simulation studies of causal effect estimation with complex longitudinal data. J. Stat. Softw. 81, 2 (2017)\label{refend:REF366}
% REF367; 原书分章页 21; RP21-L046,RP21-L047
\item\label{refentry:REF367}O. Sofrygin, M.J. van der Laan, tmlenet: targeted maximum likelihood estimation for network data (2015)\label{refend:REF367}
% REF368; 原书分章页 22; RP22-L003,RP22-L004,RP22-L005
\item\label{refentry:REF368}O. Sofrygin, M.J. van der Laan, Semi-parametric estimation and inference for the mean outcome of the single time-point intervention in a causally connected population. J. Causal Inference \textbf{5}(1), 20160003 (2017)\label{refend:REF368}
% REF369; 原书分章页 22; RP22-L006,RP22-L007
\item\label{refentry:REF369}A. Soshnikov, Gaussian limit for determinantal random point fields. Ann. Probab. \textbf{30}(1), 171--187 (2000)\label{refend:REF369}
% REF370; 原书分章页 22; RP22-L008,RP22-L009
\item\label{refentry:REF370}K. Stanley, Design of randomized controlled trials. Circulation \textbf{115}, 1164--1169 (2007)\label{refend:REF370}
% REF371; 原书分章页 22; RP22-L010,RP22-L011,RP22-L012
\item\label{refentry:REF371}R.J.C.M. Starmans, Models, inference, and truth: probabilistic reasoning in the information era, in \textit{Targeted} \textit{Learning:} \textit{Causal} \textit{Inference} \textit{for} \textit{Observational} \textit{and} \textit{Experimental} \textit{Data}, ed. by M. van der Laan, S. Rose (Springer, Berlin, 2011)\label{refend:REF371}
% REF372; 原书分章页 22; RP22-L013,RP22-L014,RP22-L015
\item\label{refentry:REF372}R.J.C.M. Starmans, The reality behind the model and the cracks in the mirror of nature (in Dutch), in \textit{Filosofie,} \textit{Tweemaandelijks} \textit{Vlaams-Nederlands} \textit{Tijdschrift,} \textit{jaargang}, vol. 21 (Garant Publishers, Antwerpen, Apeldoorn, 2011a)\label{refend:REF372}
% REF373; 原书分章页 22; RP22-L016,RP22-L017,RP22-L018
\item\label{refentry:REF373}R.J.C.M. Starmans, Ethics and statistics; the progress of a laborious dialogue (in Dutch), in \textit{Filosofie,} \textit{Tweemaandelijks} \textit{Vlaams-Nederlands} \textit{Tijdschrift,} \textit{jaargang}, vol. 22 (Garant Publishers, Antwerpen, Apeldoorn, 2012a)\label{refend:REF373}
% REF374; 原书分章页 22; RP22-L019,RP22-L020
\item\label{refentry:REF374}R.J.C.M. Starmans, Statistics, discomfort and the human dimension (in Dutch), in \textit{STAtOR}, vol. 13 (2012b)\label{refend:REF374}
% REF375; 原书分章页 22; RP22-L021,RP22-L022
\item\label{refentry:REF375}R.J.C.M. Starmans, The world of values; statistics, evolution and ethics (in Dutch), in \textit{STAtOR}, vol. 13 (2012c)\label{refend:REF375}
% REF376; 原书分章页 22; RP22-L023,RP22-L024,RP22-L025
\item\label{refentry:REF376}R.J.C.M. Starmans, Idols and ideals; francis bacon, induction and the hypothetico-deductive model (in Dutch). in \textit{Filosofie,} \textit{Tweemaandelijks} \textit{Vlaams-Nederlands} \textit{Tijdschrift,} \textit{jaargang}, vol. 23 (Garant Publishers, Antwerpen, Apeldoorn, 2013)\label{refend:REF376}
% REF377; 原书分章页 22; RP22-L026,RP22-L027,RP22-L028,RP22-L029
\item\label{refentry:REF377}R.J.C.M. Starmans, Between hobbes and turing; george boole and the laws of thinking (in Dutch), in \textit{Filosofie,} \textit{Tweemaandelijks} \textit{Vlaams-Nederlands} \textit{Tijdschrift,} \textit{jaargang}, vol. 25 (Garant Publishers, Antwerpen, Apeldoorn, 2015a)\label{refend:REF377}
% REF378; 原书分章页 22; RP22-L030,RP22-L031
\item\label{refentry:REF378}R.J.C.M. Starmans, With google toward the automatic statistician (in Dutch), in \textit{STAtOR}, vol. 16 ( 2015b)\label{refend:REF378}
% REF379; 原书分章页 22; RP22-L032,RP22-L033,RP22-L034
\item\label{refentry:REF379}R.J.C.M. Starmans, Shannon; information, entropy and the probabilistic worldview (in Dutch), in \textit{Filosofie} \textit{Tweemaandelijks} \textit{Vlaams-Nederlands} \textit{Tijdschrift,} \textit{jaargang}, vol. 26 (Garant Publishers, Antwerpen, Apeldoorn, 2016a)\label{refend:REF379}
% REF380; 原书分章页 22; RP22-L035,RP22-L036,RP22-L037,RP22-L038
\item\label{refentry:REF380}R.J.C.M. Starmans, The advent of data science - some considerations on the unreasonable effectiveness of data, in \textit{Handbook} \textit{of} \textit{Big} \textit{Data} \textit{-} \textit{Handbooks} \textit{of} \textit{Modern} \textit{Statistical} \textit{Methods}, ed. by P. Buhlmann, P. Drineas, M. Kane, M.J. van der Laan (Chapman \& Hall/CRC, New York, 2016b)\label{refend:REF380}
% REF381; 原书分章页 22; RP22-L039,RP22-L040
\item\label{refentry:REF381}R.J.C.M. Starmans, From heraclitus to shannon: the velvet revolution of data in context and flux (in Dutch), in \textit{STAtOR}, vol. 18 (2017a)\label{refend:REF381}
% REF382; 原书分章页 22; RP22-L041,RP22-L042,RP22-L043
\item\label{refentry:REF382}R.J.C.M. Starmans, The end of theory or the unreasonableness of data (in Dutch), in \textit{Filosofie,} \textit{Tweemaandelijks} \textit{Vlaams-Nederlands} \textit{Tijdschrift,} \textit{jaargang}, vol. 27 (Garant Publishers, Antwerpen, Apeldoorn, 2017b), p. 2\label{refend:REF382}
% REF383; 原书分章页 22; RP22-L044,RP22-L045,RP22-L046
\item\label{refentry:REF383}R.J.C.M. Starmans, The new house of salomon: Peter galison and the empirical tradition (in Dutch), in \textit{Filosofie,} \textit{Tweemaandelijks} \textit{Vlaams-Nederlands} \textit{Tijdschrift,} \textit{jaargang}, vol. 27 (Garant Publishers, Antwerpen, Apeldoorn, 2017c), p. 4\label{refend:REF383}
% REF384; 原书分章页 23; RP23-L003,RP23-L004,RP23-L005
\item\label{refentry:REF384}R.J.C.M. Starmans, The tryptych of the Bayesian paradigm: confirmation, inference and algoritmics, in \textit{Filosofie,} \textit{Tweemaandelijks} \textit{Vlaams-Nederlands} \textit{Tijdschrift,} \textit{jaargang}, vol. 27 (Garant Publishers, Antwerpen, Apeldoorn, 2017d)\label{refend:REF384}
% REF385; 原书分章页 23; RP23-L006,RP23-L007
\item\label{refentry:REF385}C. Steglich, T.A.B. Snijders, M. Pearson, Dynamic networks and behavior: separating selection from influence. Sociol. Methodol. \textbf{40}(1), 329--393 (2010)\label{refend:REF385}
% REF386; 原书分章页 23; RP23-L008,RP23-L009
\item\label{refentry:REF386}S. Stigler, \textit{The} \textit{History} \textit{of} \textit{Statistics:} \textit{The} \textit{Measurement} \textit{of} \textit{Uncertainty} \textit{Before} \textit{1900} (Harvard University Press, Cambridge, MA, 1986)\label{refend:REF386}
% REF387; 原书分章页 23; RP23-L010,RP23-L011
\item\label{refentry:REF387}S. Stigler, \textit{The} \textit{History} \textit{of} \textit{Statistical} \textit{Concepts} \textit{and} \textit{Methods} (Harvard University Press, Cambridge, MA, 1999)\label{refend:REF387}
% REF388; 原书分章页 23; RP23-L012,RP23-L013
\item\label{refentry:REF388}S. Stigler, \textit{The} \textit{Seven} \textit{Pillars} \textit{of} \textit{Statistical} \textit{Wisdom} (Harvard University Press, Cambridge, MA, 2016)\label{refend:REF388}
% REF389; 原书分章页 23; RP23-L014,RP23-L015,RP23-L016
\item\label{refentry:REF389}O.M. Stitelman, V. De Gruttola, M.J. van der Laan, A general implementation of TMLE for longitudinal data applied to causal inference in survival analysis. Int. J. Biostat. \textbf{8}(1), 1--37 (2012)\label{refend:REF389}
% REF390; 原书分章页 23; RP23-L017,RP23-L018
\item\label{refentry:REF390}O.M. Stitelman, M.J. van der Laan, Collaborative targeted maximum likelihood for time-to-event data. Int. J. Biostat. \textbf{6}(1), Article 21 (2010)\label{refend:REF390}
% REF391; 原书分章页 23; RP23-L019,RP23-L020
\item\label{refentry:REF391}O.M. Stitelman, M.J. van der Laan. Targeted maximum likelihood estimation of effect modification parameters in survival analysis. Int. J. Biostat. \textbf{7}(1), 1--34 (2011)\label{refend:REF391}
% REF392; 原书分章页 23; RP23-L021,RP23-L022,RP23-L023,RP23-L024
\item\label{refentry:REF392}O.M. Stitelman, V. De Gruttola, C.W. Wester, M.J. van der Laan, Rcts with time-to-event outcomes and effect modification parameters, in \textit{Targeted} \textit{Learning:} \textit{Causal} \textit{Inference} \textit{for} \textit{Observational} \textit{and} \textit{Experimental} \textit{Data}, ed. by M. J. van der Laan, S. Rose (Springer, Berlin, 2011)\label{refend:REF392}
% REF393; 原书分章页 23; RP23-L025,RP23-L026
\item\label{refentry:REF393}C.A. Struthers, J.D. Kalbfleisch, Misspecified proportional hazard models. Biometrika \textbf{73}(2), 363--369 (1986)\label{refend:REF393}
% REF394; 原书分章页 23; RP23-L027,RP23-L028,RP23-L029
\item\label{refentry:REF394}E.A. Stuart, S.R. Cole, C.P. Bradshaw, P.J. Leaf, The use of propensity scores to assess the generalizability of results from randomized trials. J. R. Stat. Soc. Ser. A \textbf{174}(Part 2), 369--386 (2011)\label{refend:REF394}
% REF395; 原书分章页 23; RP23-L030,RP23-L031
\item\label{refentry:REF395}J. Tacq, Causality in qualitative and quantitative research. Qual. Quant. \textbf{45}(2), 263--291 (2011)\label{refend:REF395}
% REF396; 原书分章页 23; RP23-L032,RP23-L033,RP23-L034
\item\label{refentry:REF396}I. Tager, M. Hollenberg, W. Satariano, Self-reported leisure-time physical activity and measures of cardiorespiratory fitness in an elderly population. Am. J. Epidemiol. \textbf{147}, 921--931 (1998)\label{refend:REF396}
% REF397; 原书分章页 23; RP23-L035,RP23-L036,RP23-L037
\item\label{refentry:REF397}E.J. Tchetgen Tchetgen, I. Shpitser, Semiparametric theory for causal mediation analysis: efficiency bounds, multiple robustness, and sensitivity analysis. Technical report 130, Biostatistics, Harvard University, June (2011a)\label{refend:REF397}
% REF398; 原书分章页 23; RP23-L038,RP23-L039,RP23-L040
\item\label{refentry:REF398}E.J. Tchetgen Tchetgen, I. Shpitser, Semiparametric estimation of models for natural direct and indirect effects. Technical Report 129, Biostatistics, Harvard University, June (2011b)\label{refend:REF398}
% REF399; 原书分章页 23; RP23-L041,RP23-L042
\item\label{refentry:REF399}E.J. Tchetgen Tchetgen, T.J. VanderWeele. On causal inference in the presence of interference. Stat. Methods Med. Res. \textbf{21}(1), 55--75 (2012)\label{refend:REF399}
% REF400; 原书分章页 23; RP23-L043,RP23-L044
\item\label{refentry:REF400}P. Thall, H. Sung, E. Estey, Selecting therapeutic strategies based on efficacy and death in multicourse clinical trials. J. Am. Stat. Assoc. \textbf{39}, 29--39 (2002)\label{refend:REF400}
% REF401; 原书分章页 23,24; RP23-L045,RP23-L046,RP24-L003,RP24-L004
\item\label{refentry:REF401}The Diabetes Control and Complications Trial Research Group, The effect of intensive treatment of diabetes on the development and progression of long-term complications in insulin-dependent diabetes mellitus. N. Engl. J. Med. \textbf{329}, 977--86 (1993)\label{refend:REF401}
% REF402; 原书分章页 24; RP24-L005,RP24-L006,RP24-L007
\item\label{refentry:REF402}M. Toftager, L.B. Christiansen, P.L. Kristensen, J. Troelsen, Space for physical activity-a multicomponent intervention study: study design and baseline findings from a cluster randomized controlled trial. BMC Public Health \textbf{11}, 777 (2011)\label{refend:REF402}
% REF403; 原书分章页 24; RP24-L008,RP24-L009
\item\label{refentry:REF403}P. Toulis, E. Kao, Estimation of causal peer influence effects, in \textit{Proceedings} \textit{of} \textit{The} \textit{30th} \textit{International} \textit{Conference} \textit{on} \textit{Machine} \textit{Learning} (2013), pp. 1489--1497\label{refend:REF403}
% REF404; 原书分章页 24; RP24-L010,RP24-L011
\item\label{refentry:REF404}A.A. Tsiatis, \textit{Semiparametric} \textit{Theory} \textit{and} \textit{Missing} \textit{Data}. (Springer, Berlin, Heidelberg, New York, 2006)\label{refend:REF404}
% REF405; 原书分章页 24; RP24-L012,RP24-L013,RP24-L014
\item\label{refentry:REF405}A.A. Tsiatis, M. Davidian, M. Zhang, X. Lu, Covariate adjustment for two-sample treatment comparisons in randomized clinical trials: a principled yet flexible approach. Stat. Med. \textbf{27}, 4658--4677 (2008)\label{refend:REF405}
% REF406; 原书分章页 24; RP24-L015,RP24-L016,RP24-L017
\item\label{refentry:REF406}C. Tuglus, M.J. van der Laan, Targeted methods for biomarker discovery, in \textit{Targeted} \textit{Learning:} \textit{Causal} \textit{Inference} \textit{for} \textit{Observational} \textit{and} \textit{Experimental} \textit{Data}. ed. by M.J. van der Laan, S. Rose (Springer, Berlin, 2011)\label{refend:REF406}
% REF407; 原书分章页 24; RP24-L018,RP24-L019,RP24-L020
\item\label{refentry:REF407}UK Prospective Diabetes Study (UKPDS) Group, Effect of intensive blood-glucose control with metformin on complications in overweight patients with type 2 diabetes (UKPDS 34). Lancet \textbf{352}, 854--865 (1998)\label{refend:REF407}
% REF408; 原书分章页 24; RP24-L021,RP24-L022,RP24-L023
\item\label{refentry:REF408}M.J. van der Laan, Causal effect models for intention to treat and realistic individualized treatment rules. Technical Report, Division of Biostatistics, University of California, Berkeley (2006a)\label{refend:REF408}
% REF409; 原书分章页 24; RP24-L024,RP24-L025
\item\label{refentry:REF409}M.J. van der Laan, Statistical inference for variable importance. Int. J. Biostat. \textbf{2}(1), Article 2 (2006b)\label{refend:REF409}
% REF410; 原书分章页 24; RP24-L026,RP24-L027
\item\label{refentry:REF410}M.J. van der Laan, Estimation based on case-control designs with known prevalence probability. Int. J. Biostat. \textbf{4}(1), Article 17 (2008a)\label{refend:REF410}
% REF411; 原书分章页 24; RP24-L028,RP24-L029,RP24-L030
\item\label{refentry:REF411}M.J. van der Laan, The construction and analysis of adaptive group sequential designs. Technical Report 232, Division of Biostatistics, University of California, Berkeley (2008b)\label{refend:REF411}
% REF412; 原书分章页 24; RP24-L031,RP24-L032
\item\label{refentry:REF412}M.J. van der Laan, Targeted maximum likelihood based causal inference: Part I. Int. J. Biostat. \textbf{6}(2), Article 2 (2010a)\label{refend:REF412}
% REF413; 原书分章页 24; RP24-L033,RP24-L034
\item\label{refentry:REF413}M.J. van der Laan, Targeted maximum likelihood based causal inference: Part II. Int. J. Biostat. \textbf{6}(2), Article 3 (2010b)\label{refend:REF413}
% REF414; 原书分章页 24; RP24-L035,RP24-L036,RP24-L037
\item\label{refentry:REF414}M.J. van der Laan, Estimation of causal effects of community-based interventions. Technical Report 268, Division of Biostatistics, University of California, Berkeley (2010c)\label{refend:REF414}
% REF415; 原书分章页 24; RP24-L038,RP24-L039
\item\label{refentry:REF415}M.J. van der Laan, Causal inference for networks. Technical Report, Division of Biostatistics, University of California, Berkeley (2012)\label{refend:REF415}
% REF416; 原书分章页 24; RP24-L040,RP24-L041
\item\label{refentry:REF416}M.J. van der Laan, Causal inference for a population of causally connected units. J. Causal Inference \textbf{2}(1), 13--74 (2014a)\label{refend:REF416}
% REF417; 原书分章页 24; RP24-L042,RP24-L043
\item\label{refentry:REF417}M.J. van der Laan, Targeted estimation of nuisance parameters to obtain valid statistical inference. Int. J. Biostat. \textbf{10}(1), 29--57 (2014b)\label{refend:REF417}
% REF418; 原书分章页 24; RP24-L044,RP24-L045
\item\label{refentry:REF418}M.J. van der Laan, A generally efficient targeted minimum loss based estimator. Int. J. Biostat. \textbf{13}(2), 1106--1118 (2017)\label{refend:REF418}
% REF419; 原书分章页 24,25; RP24-L046,RP24-L047,RP25-L003,RP25-L004
\item\label{refentry:REF419}M.J. van der Laan, S. Dudoit, Unified cross-validation methodology for selection among estimators and a general cross-validated adaptive epsilon-net estimator: finite sample oracle inequalities and examples. Technical Report, Division of Biostatistics, University of California, Berkeley (2003)\label{refend:REF419}
% REF420; 原书分章页 25; RP25-L005,RP25-L006
\item\label{refentry:REF420}M.J. van der Laan, S. Gruber, Collaborative double robust penalized targeted maximum likelihood estimation. Int. J. Biostat. \textbf{6}(1), Article 17 (2010)\label{refend:REF420}
% REF421; 原书分章页 25; RP25-L007,RP25-L008
\item\label{refentry:REF421}M.J. van der Laan, S. Gruber, Targeted minimum loss based estimation of causal effects of multiple time point interventions. Int. J. Biostat. \textbf{8}(1), Article 9 (2012)\label{refend:REF421}
% REF422; 原书分章页 25; RP25-L009,RP25-L010,RP25-L011
\item\label{refentry:REF422}M.J. van der Laan, S. Gruber, One-step targeted minimum loss-based estimation based on universal least favorable one-dimensional submodels. Int. J. Biostat. \textbf{12}(1), 351--378 (2016)\label{refend:REF422}
% REF423; 原书分章页 25; RP25-L012,RP25-L013
\item\label{refentry:REF423}M.J. van der Laan, S. Lendle, Online targeted learning. Technical Report, Division of Biostatistics, University of California, Berkeley (2014)\label{refend:REF423}
% REF424; 原书分章页 25; RP25-L014,RP25-L015,RP25-L016
\item\label{refentry:REF424}M.J. van der Laan, A.R. Luedtke, Targeted learning of an optimal dynamic treatment, and statistical inference for its mean outcome. Technical Report, Division of Biostatistics, University of California, Berkeley\label{refend:REF424}
% REF425; 原书分章页 25; RP25-L017,RP25-L018
\item\label{refentry:REF425}M.J. van der Laan, A.R. Luedtke, Targeted learning of the mean outcome under an optimal dynamic treatment rule. J. Causal Inference \textbf{3}(1), 61--95 (2015)\label{refend:REF425}
% REF426; 原书分章页 25; RP25-L019,RP25-L020
\item\label{refentry:REF426}M.J. van der Laan, M.L. Petersen, Causal effect models for realistic individualized treatment and intention to treat rules. Int. J. Biostat. \textbf{3}(1), Article 3 (2007)\label{refend:REF426}
% REF427; 原书分章页 25; RP25-L021,RP25-L022
\item\label{refentry:REF427}M.J. van der Laan, M.L. Petersen, Direct effect models. Int. J. Biostat. \textbf{4}(1), Article 23 (2008)\label{refend:REF427}
% REF428; 原书分章页 25; RP25-L023,RP25-L024
\item\label{refentry:REF428}M.J. van der Laan, K.S. Pollard, Hybrid clustering of gene expression data with visualization and the bootstrap. J. Stat. Plann. Inference \textbf{117}, 275--303 (2003)\label{refend:REF428}
% REF429; 原书分章页 25; RP25-L025,RP25-L026
\item\label{refentry:REF429}M.J. van der Laan, E.C. Polley, A.E. Hubbard, Super learner. Stat. Appl. Genet. Mol. \textbf{6}(1), Article 25 (2007)\label{refend:REF429}
% REF430; 原书分章页 25; RP25-L027,RP25-L028
\item\label{refentry:REF430}M.J. van der Laan, J.M. Robins, \textit{Unified} \textit{Methods} \textit{for} \textit{Censored} \textit{Longitudinal} \textit{Data} \textit{and} \textit{Causality} (Springer, Berlin Heidelberg New York, 2003)\label{refend:REF430}
% REF431; 原书分章页 25; RP25-L029,RP25-L030
\item\label{refentry:REF431}M.J. van der Laan, S. Rose, \textit{Targeted} \textit{Learning:} \textit{Causal} \textit{Inference} \textit{for} \textit{Observational} \textit{and} \textit{Experimental} \textit{Data} (Springer, Berlin, Heidelberg, New York, 2011)\label{refend:REF431}
% REF432; 原书分章页 25; RP25-L031,RP25-L032
\item\label{refentry:REF432}M.J. van der Laan, D.B. Rubin, Targeted maximum likelihood learning. Int. J. Biostat. \textbf{2}(1), Article 11 (2006)\label{refend:REF432}
% REF433; 原书分章页 25; RP25-L033,RP25-L034,RP25-L035
\item\label{refentry:REF433}M.J. van der Laan, R.J.C.M. Starmans, Entering the era of data science: targeted learning and the integration of statistics and computational data analysis. Adv. Stat. 2014, 502678 (2014)\label{refend:REF433}
% REF434; 原书分章页 25; RP25-L036,RP25-L037
\item\label{refentry:REF434}M.J. van der Laan, S. Dudoit, S. Keleş, Asymptotic optimality of likelihood-based cross-validation. Stat. Appl. Genet. Mol. \textbf{3}(1), Article 4 (2004)\label{refend:REF434}
% REF435; 原书分章页 25; RP25-L038,RP25-L039
\item\label{refentry:REF435}M.J. van der Laan, S. Dudoit, A.W. van der Vaart. The cross-validated adaptive epsilon-net estimator. Stat. Decis. \textbf{24}(3), 373--395 (2006)\label{refend:REF435}
% REF436; 原书分章页 25; RP25-L040,RP25-L041
\item\label{refentry:REF436}M.J. van der Laan, L.B. Balzer, M.L. Petersen, Adaptive matching in randomized trials and observational studies. J. Stat. Res. \textbf{46}(2), 113--156 (2013a)\label{refend:REF436}
% REF437; 原书分章页 25; RP25-L042,RP25-L043,RP25-L044
\item\label{refentry:REF437}M.J. van der Laan, M. Petersen, W. Zheng, Estimating the effect of a community-based intervention with two communities. J. Causal Inference \textbf{1}(1), 83--106 (2013b)\label{refend:REF437}
% REF438; 原书分章页 25,26; RP25-L045,RP25-L046,RP26-L003,RP26-L004
\item\label{refentry:REF438}M.J. van der Laan, A.R. Luedtke, I. Díaz, Discussion of identification, estimation and approximation of risk under interventions that depend on the natural value of treatment using observational data, by Jessica Young, Miguel Hernán, and James Robins. Epidemiol Methods \textbf{3}(1), 21--31 (2014)\label{refend:REF438}
% REF439; 原书分章页 26; RP26-L005,RP26-L006,RP26-L007
\item\label{refentry:REF439}M.J. van der Laan, M. Carone, A.R. Luedtke, Computerizing efficient estimation of a pathwise differentiable target parameter. Technical Report, Division of Biostatistics, University of California, Berkeley (2015)\label{refend:REF439}
% REF440; 原书分章页 26; RP26-L008
\item\label{refentry:REF440}A.W. van der Vaart, \textit{Asymptotic} \textit{Statistics} (Cambridge, New York, 1998)\label{refend:REF440}
% REF441; 原书分章页 26; RP26-L009,RP26-L010
\item\label{refentry:REF441}A.W. van der Vaart, Higher order tangent spaces and influence functions. Stat. Sci. \textbf{29}(4), 679--686 (2014)\label{refend:REF441}
% REF442; 原书分章页 26; RP26-L011,RP26-L012
\item\label{refentry:REF442}A.W. van der Vaart, J.A. Wellner, \textit{Weak} \textit{Convergence} \textit{and} \textit{Empirical} \textit{Processes} (Springer, Berlin, Heidelberg, New York, 1996)\label{refend:REF442}
% REF443; 原书分章页 26; RP26-L013,RP26-L014
\item\label{refentry:REF443}A.W. van der Vaart, J.A. Wellner, A local maximal inequality under uniform entropy. Electron. J. Stat. \textbf{5}, 192--203 (2011)\label{refend:REF443}
% REF444; 原书分章页 26; RP26-L015,RP26-L016
\item\label{refentry:REF444}A.W. van der Vaart, S. Dudoit, M.J. van der Laan, Oracle inequalities for multi-fold cross-validation. Stat. Decis. \textbf{24}(3), 351--371 (2006)\label{refend:REF444}
% REF445; 原书分章页 26; RP26-L017,RP26-L018
\item\label{refentry:REF445}R. van Handel, On the minimal penalty for Markov order estimation. Probab. Theory Relat. Fields \textbf{150}, 709--738 (2009)\label{refend:REF445}
% REF446; 原书分章页 26; RP26-L019,RP26-L020
\item\label{refentry:REF446}T.J. VanderWeele, Marginal structural models for the estimation of direct and indirect effects. Epidemiology \textbf{20}, 18--26 (2009)\label{refend:REF446}
% REF447; 原书分章页 26; RP26-L021,RP26-L022
\item\label{refentry:REF447}T.J. VanderWeele, Bias formulas for sensitivity analysis for direct and indirect effects. Epidemiology \textbf{21}(4), 540 (2010)\label{refend:REF447}
% REF448; 原书分章页 26; RP26-L023,RP26-L024
\item\label{refentry:REF448}T.J VanderWeele, Sensitivity analysis for contagion effects in social networks. Sociol. Methods Res. \textbf{40}(2), 240--255 (2011)\label{refend:REF448}
% REF449; 原书分章页 26; RP26-L025,RP26-L026
\item\label{refentry:REF449}T.J. VanderWeele, Inference for influence over multiple degrees of separation on a social network. Stat. Med. \textbf{32}(4), 591--596 (2013)\label{refend:REF449}
% REF450; 原书分章页 26; RP26-L027,RP26-L028
\item\label{refentry:REF450}T.J. VanderWeele, W. An, Social networks and causal inference, in \textit{Handbook} \textit{of} \textit{Causal} \textit{Analysis} \textit{for} \textit{Social} \textit{Research} (Springer, Berlin, 2013), pp. 353--374\label{refend:REF450}
% REF451; 原书分章页 26; RP26-L029,RP26-L030,RP26-L031
\item\label{refentry:REF451}T.J. VanderWeele, O.A. Arah, Unmeasured confounding for general outcomes, treatments, and confounders: bias formulas for sensitivity analysis. Epidemiology \textbf{22}(1), 42 (2011)\label{refend:REF451}
% REF452; 原书分章页 26; RP26-L032,RP26-L033
\item\label{refentry:REF452}T.J. VanderWeele, M.A. Hernán, Causal inference under multiple versions of treatment. J. Causal Inference \textbf{1}(1), 1--20 (2013)\label{refend:REF452}
% REF453; 原书分章页 26; RP26-L034,RP26-L035
\item\label{refentry:REF453}T.J. VanderWeele, E.J. Tchetgen Tchetgen, Mediation analysis with time-varying exposures and mediators. J. R. Stat. Soc. Ser. B \textbf{79}(3), 917--938 (2017)\label{refend:REF453}
% REF454; 原书分章页 26; RP26-L036,RP26-L037
\item\label{refentry:REF454}T.J. VanderWeele, B. Mukherjee, J. Chen, Sensitivity analysis for interactions under unmeasured confounding. Stat. Med. \textbf{31}(22), 2552--2564 (2012a)\label{refend:REF454}
% REF455; 原书分章页 26; RP26-L038,RP26-L039,RP26-L040
\item\label{refentry:REF455}T.J. VanderWeele, J.P. Vandenbrouke, E.J. Tchetgen Tchetgen, J.M. Robins, A mapping between interactions and interference: implications for vaccine trials. Epidemiology \textbf{23}(3), 285--292 (2012b)\label{refend:REF455}
% REF456; 原书分章页 26; RP26-L041,RP26-L042,RP26-L043
\item\label{refentry:REF456}T.J. VanderWeele, E.L. Ogburn, E.J. Tchetgen Tchetgen, Why and when ``flawed" social network analyses still yield valid tests of no contagion. Stat. Polit. Policy \textbf{3}(1), 2151--2160 (2012c)\label{refend:REF456}
% REF457; 原书分章页 26; RP26-L044,RP26-L045,RP26-L046
\item\label{refentry:REF457}T.J. VanderWeele, S. Vansteelandt, J.M. Robins, Effect decomposition in the presence of an exposure-induced mediator-outcome confounder. Epidemiology \textbf{25}(2), 300--306 (2014a)\label{refend:REF457}
% REF458; 原书分章页 27; RP27-L003,RP27-L004
\item\label{refentry:REF458}T.J. VanderWeele, E.J. Tchetgen Tchetgen, M.E. Halloran, Interference and sensitivity analysis. Stat. Sci. \textbf{29}(4), 687--706 (2014b)\label{refend:REF458}
% REF459; 原书分章页 27; RP27-L005,RP27-L006,RP27-L007
\item\label{refentry:REF459}S. Vansteelandt, E. Goetghebeur, M.G. Kenward, G. Molenberghs, Ignorance and uncertainty regions as inferential tools in a sensitivity analysis. Stat. Sin. \textbf{16}(3), 953--979 (2006)\label{refend:REF459}
% REF460; 原书分章页 27; RP27-L008,RP27-L009
\item\label{refentry:REF460}W.N. Venables, B.D. Ripley, \textit{Modern} \textit{Applied} \textit{Statistics} \textit{with} \textit{S}, 4th edn. (Springer, Berlin, Heidelberg, New York, 2002)\label{refend:REF460}
% REF461; 原书分章页 27; RP27-L010,RP27-L011,RP27-L012
\item\label{refentry:REF461}T.M. Vogt, J. Elston-Lafata, D. Tolsma, S.M. Greene, The role of research in integrated healthcare systems: the HMO Research Network. Am. J. Manag. Care \textbf{10}(9), 643--648 (2004)\label{refend:REF461}
% REF462; 原书分章页 27; RP27-L013,RP27-L014
\item\label{refentry:REF462}E. Wagner, B. Austin, C. Davis, M. Hindmarsh, J. Schaefer, A. Bonomi, Improving chronic illness care: translating evidence into action. Health Aff. \textbf{20}, 64--78 (2001)\label{refend:REF462}
% REF463; 原书分章页 27; RP27-L015,RP27-L016
\item\label{refentry:REF463}D. Walker, L. Muchnik, Design of randomized experiments in networks. Proc. IEEE \textbf{102}(12), 1940--1951 (2014)\label{refend:REF463}
% REF464; 原书分章页 27; RP27-L017,RP27-L018
\item\label{refentry:REF464}H. Wang, M.J. van der Laan, Dimension reduction with gene expression data using targeted variable importance measurement. BMC Bioinf. \textbf{12}(1), 312 (2011)\label{refend:REF464}
% REF465; 原书分章页 27; RP27-L019,RP27-L020,RP27-L021
\item\label{refentry:REF465}H. Wang, S. Rose, M.J. van der Laan, Finding quantitative trait loci genes with collaborative targeted maximum likelihood learning. Stat. Probab. Lett. \textbf{81}(7), 792--796 (2011a)\label{refend:REF465}
% REF466; 原书分章页 27; RP27-L022,RP27-L023,RP27-L024
\item\label{refentry:REF466}H. Wang, S. Rose, M.J. van der Laan. Finding quantitative trait loci genes, in \textit{Targeted} \textit{Learning:} \textit{Causal} \textit{Inference} \textit{for} \textit{Observational} \textit{and} \textit{Experimental} \textit{Data}, ed. by M.J. van der Laan, S. Rose (Springer, Berlin Heidelberg, New York, 2011b)\label{refend:REF466}
% REF467; 原书分章页 27; RP27-L025,RP27-L026
\item\label{refentry:REF467}H. Wang, Z. Zhang, S. Rose, M.J. van der Laan, A novel targeted learning methods for quantitative trait Loci mapping. Genetics \textbf{198}(4), 1369--1376 (2014)\label{refend:REF467}
% REF468; 原书分章页 27; RP27-L027,RP27-L028
\item\label{refentry:REF468}G.S. Watson, Smooth regression analysis. Sankhyā Indian J. Stat. Ser. A 359--372 (1964)\label{refend:REF468}
% REF469; 原书分章页 27; RP27-L029,RP27-L030,RP27-L031
\item\label{refentry:REF469}L. Watson, R. Small, S. Brown, W. Dawson, J. Lumley, Mounting a community-randomized trial: sample size, matching, selection, and randomization issues in PRISM. Control. Clin. Trials \textbf{25}(3), 235--250 (2004)\label{refend:REF469}
% REF470; 原书分章页 27; RP27-L032,RP27-L033
\item\label{refentry:REF470}S. Weinberg, \textit{Dreams} \textit{of} \textit{a} \textit{Final} \textit{Theory:} \textit{The} \textit{Scientist's} \textit{Search} \textit{for} \textit{the} \textit{Ultimate} \textit{Laws} \textit{of} \textit{Nature} (Random House Inc., New York, 1993)\label{refend:REF470}
% REF471; 原书分章页 27; RP27-L034,RP27-L035
\item\label{refentry:REF471}D. Wied, R. Weißbach, Consistency of the kernel density estimator: a survey. Stat. Pap. \textbf{53}(1), 1--21 (2012)\label{refend:REF471}
% REF472; 原书分章页 27; RP27-L036,RP27-L037
\item\label{refentry:REF472}R.J. Wieringa, Design Science Methodology for Information Systems and Software Engineering (Springer, New York, 2014)\label{refend:REF472}
% REF473; 原书分章页 27; RP27-L038,RP27-L039,RP27-L040
\item\label{refentry:REF473}J. Williamson, Probabilistic theories of causality, in \textit{The} \textit{Oxford} \textit{Handbook} \textit{of} \textit{Causation}, ed. by H. Beebee, C. Hitchcock, P. Menzies (Oxford University Press, Oxford, 2009), pp. 185--212\label{refend:REF473}
% REF474; 原书分章页 27; RP27-L041,RP27-L042,RP27-L043
\item\label{refentry:REF474}P. Wilson, R.B. D'Agostino, D. Levy, A.M. Belanger, H. Silbershatz, W.B. Kannel, Prediction of coronary heart disease using risk factor categories. Circulation \textbf{97}(18), 1837--1847 (1998)\label{refend:REF474}
% REF475; 原书分章页 27; RP27-L044,RP27-L045
\item\label{refentry:REF475}T Woutersen, A simple way to calculate confidence intervals for partially identified parameters. Technical Report, Johns Hopkins University (2006)\label{refend:REF475}
% REF476; 原书分章页 27; RP27-L046,RP27-L047
\item\label{refentry:REF476}W. Xu, Towards optimal one pass large scale learning with averaged stochastic gradient descent. ArXiv e-prints, December (2011)\label{refend:REF476}
% REF477; 原书分章页 28; RP28-L003,RP28-L004,RP28-L005
\item\label{refentry:REF477}J.G. Young, M.A. Hernán, J.M. Robins, Identification, estimation and approximation of risk under interventions that depend on the natural value of treatment using observational data. Epidemiol. Methods \textbf{3}(1), 1--19 (2014)\label{refend:REF477}
% REF478; 原书分章页 28; RP28-L006,RP28-L007,RP28-L008
\item\label{refentry:REF478}S. Yuan, H.H. Zhang, M. Davidian, Variable selection for covariate-adjusted semi-parametric inference in randomized clinical trials. Stat. Med. \textbf{31}, 3789--3804 (2012)\label{refend:REF478}
% REF479; 原书分章页 28; RP28-L009,RP28-L010
\item\label{refentry:REF479}M.D. Zeiler, Adadelta: an adaptive learning rate method. arXiv e-prints, December (2012)\label{refend:REF479}
% REF480; 原书分章页 28; RP28-L011,RP28-L012,RP28-L013
\item\label{refentry:REF480}K. Zhang, D.S. Small, Comment: the essential role of pair matching in cluster-randomized experiments, with application to the Mexican universal health insurance evaluation. Stat. Sci. \textbf{25}(1), 59--64 (2009)\label{refend:REF480}
% REF481; 原书分章页 28; RP28-L014,RP28-L015
\item\label{refentry:REF481}B. Zhang, A. Tsiatis, M. Davidian, M. Zhang, E. Laber, A robust method for estimating optimal treatment regimes. Biometrics \textbf{68}, 1010--1018 (2012a)\label{refend:REF481}
% REF482; 原书分章页 28; RP28-L016,RP28-L017
\item\label{refentry:REF482}B. Zhang, A. Tsiatis, M. Davidian, M. Zhang, E. Laber, Estimating optimal treatment regimes from a classification perspective. Stat \textbf{68}(1), 103--114 (2012b)\label{refend:REF482}
% REF483; 原书分章页 28; RP28-L018,RP28-L019
\item\label{refentry:REF483}M. Zhang, A.A. Tsiatis, M. Davidian, Improving efficiency of inferences in randomized clinical trials using auxiliary covariates. Biometrics \textbf{64}(3), 707--715 (2008)\label{refend:REF483}
% REF484; 原书分章页 28; RP28-L020,RP28-L021,RP28-L022
\item\label{refentry:REF484}T. Zhang, J. Wu, F. Li, B. Caffo, D. Boatman-Reich, A dynamic directional model for effective brain connectivity using electrocorticographic (ECoG) time series. J. Am. Stat. Assoc. \textbf{110}(509), 93--106 (2015)\label{refend:REF484}
% REF485; 原书分章页 28; RP28-L023,RP28-L024
\item\label{refentry:REF485}Y. Zhao, D. Zeng, A. Rush, M Kosorok, Estimating individual treatment rules using outcome weighted learning. J. Am. Stat. Assoc. \textbf{107}, 1106--1118 (2012)\label{refend:REF485}
% REF486; 原书分章页 28; RP28-L025,RP28-L026,RP28-L027
\item\label{refentry:REF486}Y. Zhao, D. Zeng, E.B. Laber, M.R. Kosorok, New statistical learning methods for estimating optimal dynamic treatment regimes. J. Am. Stat. Assoc. \textbf{110}(510), 583--598 (2015)\label{refend:REF486}
% REF487; 原书分章页 28; RP28-L028,RP28-L029,RP28-L030
\item\label{refentry:REF487}W. Zheng, M.J. van der Laan, Asymptotic theory for cross-validated targeted maximum likelihood estimation. Technical Report, Division of Biostatistics, University of California, Berkeley (2010)\label{refend:REF487}
% REF488; 原书分章页 28; RP28-L031,RP28-L032,RP28-L033
\item\label{refentry:REF488}W. Zheng, M.J. van der Laan, Causal mediation in a survival setting with time-dependent mediators. Technical Report, Division of Biostatistics, University of California, Berkeley (2012a)\label{refend:REF488}
% REF489; 原书分章页 28; RP28-L034,RP28-L035
\item\label{refentry:REF489}W. Zheng, M.J. van der Laan, Targeted maximum likelihood estimation of natural direct effects. Int. J. Biostat. \textbf{8}(1), 1--40 (2012b)\label{refend:REF489}
% REF490; 原书分章页 28; RP28-L036,RP28-L037,RP28-L038
\item\label{refentry:REF490}W. Zheng, M.J. van der Laan, Longitudinal mediation analysis with time-varying mediators and exposures, with application to survival outcomes. J. Causal Inference \textbf{5}(2), 20160006 (2017)\label{refend:REF490}
% REF491; 原书分章页 28; RP28-L039,RP28-L040,RP28-L041,RP28-L042
\item\label{refentry:REF491}W. Zheng, A. Chambaz, M.J. van der Laan, Drawing valid targeted inference when covariate-adjusted response-adaptive RCT meets data-adaptive loss-based estimation, with an application to the lasso. Technical Report, Division of Biostatistics, University of California, Berkeley (2015)\label{refend:REF491}
% REF492; 原书分章页 28; RP28-L043,RP28-L044
\item\label{refentry:REF492}M. Zinkevich, Online convex programming and generalized infinitesimal gradient ascent. Proceedings of ICML (2003)\label{refend:REF492}
\end{list}
\par\endgroup
